\chapter{Evolution Equations}
\label{chap:cz-evolution-equations}

%%\newtheorem{Definition}{Definition}[section]
%%\newtheorem{Theorem}{Theorem}[section]
%%\newtheorem{Remark}{Remark}[section]
%%\newtheorem{Lemma}{Lemma}[section]
%%\newtheorem{Proposition}{Proposition}[section]
%%\newenvironment{pf}{{\medskip\noindent \bf\em Proof.}}{{\hfill$\square$}\\}
%%$\#$}\\}

%%\numberwithin{section}{chapter}

In this chapter, we collect some basic material on Hamilton's
Ricci flow. In Section 1.1, we introduce the Ricci flow equation
of Hamilton \cite{Ha82} and examine some special solutions. The
short time existence and uniqueness theorem of the Ricci flow on a
compact manifold \cite{Ha82} (also cf. \cite{De}) is presented in
Section 1.2. Evolution equations of curvatures under the Ricci
flow \cite{Ha82, Ha86} are described in Section 1.3. In Section
1.4, we recall Shi's local derivative estimate \cite{Sh89}, which
plays an important role in the Ricci flow. Perelman's two
functionals in \cite{P1} and their monotonicity properties are
discussed in Section 1.5.



\section{The Ricci Flow}

In this section, we introduce Hamilton's Ricci flow and examine
some special solutions. The presentation follows closely
\cite{Ha82, Ha95F}.

Let $M$ be an $n$-dimensional complete Riemannian manifold with
the Riemannian metric $g_{ij}$.  The Levi-Civita connection is
given by the Christoffel symbols
$$
\Gamma^k_{ij}=\frac{1}{2}g^{kl}\left( \frac{\partial
g_{jl}}{\partial x^i}+\frac{\partial g_{il}}{\partial
x^j}-\frac{\partial g_{ij}}{\partial x^l}\right)
$$
where $g^{ij}$ is the inverse of $g_{ij}$. The summation
convention of summing over repeated indices is used here and
throughout the book. The Riemannian curvature tensor is given by
$$
R^k_{ijl}=\frac{\partial \Gamma^k_{jl}}{\partial x^i}
-\frac{\partial \Gamma^k_{il}}{\partial
x^j}+\Gamma^k_{ip}\Gamma^p_{jl}-\Gamma^k_{jp}\Gamma^p_{il}.
$$
We lower the index to the third position, so that
$$
R_{ijkl}=g_{kp}R^p_{ijl}.
$$
The curvature tensor $R_{ijkl}$ is anti-symmetric in the pairs
$i,\; j$ and $k,\; l$ and symmetric in their interchange:
$$
R_{ijkl}=-R_{jikl}=-R_{ijlk}=R_{klij}.
$$
Also the first Bianchi identity holds \be
R_{ijkl}+R_{jkil}+R_{kijl}=0.              %%\eqno(1.1.1)
\ee The Ricci tensor is the contraction
$$
R_{ik}=g^{jl}R_{ijkl},
$$
and the scalar curvature is
$$
R=g^{ij}R_{ij}.
$$

We denote the covariant derivative of a vector field
$v=v^j\frac{\partial}{\partial x^j}$ by
$$
\nabla_iv^j=\frac{\partial v^j}{\partial x^i}+\Gamma^j_{ik}v^k
$$
and of a $1$-form by
$$
\nabla_iv_j=\frac{\partial v_j}{\partial x^i}-\Gamma^k_{ij}v_k.
$$
These definitions extend uniquely to tensors so as to preserve the
product rule and contractions. For the exchange of two covariant
derivatives, we have
\begin{align}
\nabla_i\nabla_jv^l-\nabla_j\nabla_iv^l&=R^l_{ijk}v^k,\\
%\eqno(1.1.2)
\nabla_i\nabla_jv_k-\nabla_j\nabla_iv_k&=R_{ijkl}g^{lm}v_m,
%\eqno(1.1.3)
\end{align}
and similar formulas for more complicated tensors.
 The second Bianchi identity is given by
\be \nabla_mR_{ijkl}+\nabla_iR_{jmkl}+\nabla_jR_{mikl}=0.
%\eqno(1.1.4)
\ee

For any tensor $T=T^i_{jk}$ we define its length by
$$
|T^i_{jk}|^2=g_{il}g^{jm}g^{kp}T^i_{jk}T^l_{mp},
$$
and we define its Laplacian by
$$
\Delta T^i_{jk}=g^{pq}\nabla_p\nabla_qT^i_{jk},
$$
the trace of the second iterated covariant derivatives. Similar
definitions hold for more general tensors.

The {\bf Ricci flow}\index{Ricci flow} of Hamilton \cite{Ha82} is
the evolution equation \be
\frac{\partial g_{ij}}{\partial t}=-2R_{ij} %\eqno (1.1.5)
\ee for a family of Riemannian metrics $g_{ij}(t)$ on $M$. It is a
nonlinear system of second order partial differential equations on
metrics.

In order to get a feel for the Ricci flow (1.1.5) we first present
some examples of specific solutions (cf. Section 2 of
\cite{Ha95F}). \vskip 0.2cm

(1) Einstein metrics \vskip 0.1cm

A Riemannian metric $g_{ij}$ is called {\bf
Einstein}\index{Einstein! metric} if
$$
R_{ij}=\lambda g_{ij}
$$
for some constant $\lambda$. A smooth manifold $M$ with an
Einstein metric is called an {\bf Einstein
manifold}\index{Einstein! manifold}.

If the initial metric is Ricci flat, so that $R_{ij}=0$, then
clearly the metric does not change under (1.1.5). Hence any Ricci
flat metric is a stationary solution of the Ricci flow. This
happens, for example, on a flat torus or on any $K3$-surface with
a Calabi-Yau metric.

If the initial metric is Einstein with positive scalar curvature,
then the metric will shrink under the Ricci flow by a
time-dependent factor. Indeed, since the initial metric is
Einstein, we have
$$
R_{ij}(x,0)=\lambda g_{ij}(x,0), \quad \forall x\in M
$$
and some $\lambda>0$. Let
$$
g_{ij}(x,t)=\rho^2(t)g_{ij}(x,0).
$$
{}From the definition of the Ricci tensor, one sees that
$$
R_{ij}(x,t)=R_{ij}(x,0)=\lambda g_{ij}(x,0).
$$
Thus the equation (1.1.5) corresponds to
$$
\frac{\partial (\rho^2(t) g_{ij}(x,0))}{\partial t}=-2\lambda
g_{ij}(x,0).
$$
This gives the ODE \be
\frac{d\rho}{dt}=-\frac{\lambda}{\rho}, %\eqno (1.1.6)
\ee whose solution is given by
$$
\rho^2(t)=1-2\lambda t.
$$
Thus the evolving metric $g_{ij}(x,t)$ shrinks homothetically to a
point as $t\rightarrow T=1/{2\lambda}$. Note that as $t\rightarrow
T$, the scalar curvature becomes infinite like $1/(T-t)$.

By contrast, if the initial metric is an Einstein metric of
negative scalar curvature, the metric will expand homothetically
for all times. Indeed if
$$
R_{ij}(x,0)=-\lambda g_{ij}(x,0)
$$
with $\lambda >0$ and
$$
g_{ij}(x,t)=\rho^2(t)g_{ij}(x,0).
$$
Then $\rho(t)$ satisfies the ODE \be
\frac{d\rho}{dt}=\frac{\lambda}{\rho}, %\eqno (1.1.7)
\ee with the solution
$$
\rho^2(t)=1+2\lambda t.
$$
Hence the evolving metric $g_{ij}(x,t)=\rho^2(t)g_{ij}(x,0)$
exists and expands homothetically for all times, and the curvature
will fall back to zero like $-1/t$. Note that now the evolving
metric $g_{ij}(x,t)$ only goes back in time to $-1/{2\lambda}$,
when the metric explodes out of a single point in a ``big bang".
\vskip 0.2cm

(2) Ricci Solitons \vskip 0.1cm

The concept of a Ricci soliton is introduced by Hamilton
\cite{Ha88}. We will call a solution to an evolution equation
which moves under a one-parameter subgroup of the symmetry group
of the equation a {\bf steady soliton}\index{soliton!steady}. The
symmetry group of the Ricci flow contains the full diffeomorphism
group. Thus a solution to the Ricci flow (1.1.5) which moves by a
one-parameter group of diffeomorphisms $\varphi_t$ is called a
\textbf{steady Ricci soliton}\index{Ricci soliton!steady }.

If $\varphi_t$ is a one-parameter group of diffeomorphisms
generated by a vector field $V$ on $M$, then the Ricci soliton is
given by \be
g_{ij}(x,t)=\varphi^*_tg_{ij}(x,0) %\eqno (1.1.8)
\ee which implies that the Ricci term $-2 Ric$ on the RHS of
(1.1.5) is equal to the Lie derivative $\mathcal L_V g$ of the
evolving metric $g$. In particular, the initial metric
$g_{ij}(x,0)$ satisfies the following steady Ricci soliton
equation \be
2R_{ij}+g_{ik}\nabla_j V^k+g_{jk}\nabla_i V^k=0 .%\eqno (1.1.9)
\ee If the vector field $V$ is the gradient of a function $f$ then
the soliton is called a {\bf steady gradient Ricci
soliton}\index{Ricci soliton!gradient }. Thus \be
R_{ij}+\nabla_i\nabla_jf=0,
\quad \text{or}\quad Ric+{\nabla}^2f=0, %\eqno (1.1.10)
\ee is the steady gradient Ricci soliton equation.

Conversely, it is clear that a metric $g_{ij}$ satisfying (1.1.10)
generates a steady gradient Ricci soliton $g_{ij}(t)$ given by
(1.1.8). For this reason we also often call such a metric $g_{ij}$
a steady gradient Ricci soliton and do not necessarily distinguish
it with the solution $g_{ij} (t)$ it generates.

More generally, we can consider a solution to the Ricci flow
(1.1.5) which moves by diffeomorphisms and also shrinks or expands
by a (time-dependent) factor at the same time. Such a solution is
called a homothetically {\bf shrinking}\index{Ricci soliton!
shrinking} or homothetically {\bf expanding}\index{Ricci soliton!
expanding} \textbf{Ricci soliton}. The equation for a homothetic
Ricci soliton is \be 2R_{ij}+g_{ik}\nabla_j V^k+g_{jk}\nabla_i
V^k-2\lambda g_{ij}=0,
%\eqno (1.1.11)
\ee or for a homothetic gradient Ricci soliton, \be
R_{ij}+\nabla_i\nabla_jf-\lambda g_{ij}=0, %\eqno (1.1.12)
\ee where $\lambda$ is the homothetic constant. For $\lambda>0$
the soliton is shrinking, for $\lambda<0$ it is expanding. The
case $\lambda=0$ is a steady Ricci soliton, the case $V=0$ (or $f$
being a constant function) is an Einstein metric. Thus Ricci
solitons can be considered as natural extensions of Einstein
metrics. In fact, the following result states that there are no
nontrivial gradient steady or expanding Ricci solitons on any
compact manifold.

We remark that if the underlying manifold $M$ is a complex
manifold and the initial metric is K\"ahler, then it is well known
(see, e.g., \cite {Ha95, Cao85}) that the solution metric to the
Ricci flow (1.1.5) remains K\"ahler. For this reason, the Ricci
flow on a K\"ahler manifold is called the {\bf K\"ahler-Ricci
flow}\index{K\"ahler-Ricci flow}. A (steady, or shrinking, or
expanding) Ricci soliton to the K\"ahler-Ricci flow is called a
({\bf steady}, or {\bf shrinking}, or {\bf expanding} repectively)
{\bf K\"ahler-Ricci soliton}. \index{K\"ahler-Ricci soliton!
steady} \index{K\"ahler-Ricci soliton! shrinking}
\index{K\"ahler-Ricci soliton! expanding} The following result,
essentially due to Hamilton (cf. Theorem 20.1 of \cite{Ha95F}, or
Proposition 5.20 of \cite{CK04}), says there are no nontrivial
compact steady and expanding Ricci solitons.


%%%\vskip 0.3cm \noindent {\bf Proposition 1.1.1}\ \  \emph{
% CZ-BASELINE source-line:784 source-kind:proposition source-id:1.1.1
\begin{proposition}[compact steady and expanding gradient solitons are Einstein]
  \label{prop:cz-compact-steady-and-expanding-gradient-solitons-are-einstein}
  \dcref{Proposition 1.1.1}
  \group{Cao--Zhu The Ricci Flow}
  \source{cao-zhu}
On a compact n-dimensional manifold $M$, a gradient steady or
expanding Ricci soliton is necessarily an Einstein metric.
\end{proposition}

%%%\vskip 0.1cm \noindent {\bf\em Proof.}
\begin{proof}
We first prove the steady case. Our argument here follows that of Hamilton
\cite{Ha95F}.

Let $g_{ij}$ be a complete steady gradient Ricci soliton on a
manifold $M$ so that
$$
R_{ij}+\nabla_i\nabla_jf=0.
$$

Taking the trace, we get \be
R+\Delta f=0. %\eqno(1.1.13)
\ee Also, taking the covariant derivatives of the Ricci soliton
equation, we have
$$
\nabla_i\nabla_j\nabla_kf-\nabla_j\nabla_i\nabla_kf
=\nabla_jR_{ik}-\nabla_iR_{jk}.
$$
On the other hand, by using the commutating formula (1.1.3), we
otain
$$
\nabla_i\nabla_j\nabla_kf-\nabla_j\nabla_i\nabla_kf
=R_{ijkl}\nabla_lf.
$$
Thus
$$
\nabla_iR_{jk}-\nabla_jR_{ik}+R_{ijkl}\nabla_lf=0.
$$
Taking the trace on $j$ and $k$, and using the contracted second
Bianchi identity \be
\nabla_j R_{ij}=\frac {1}{2}\nabla_i R, %\eqno(1.1.14)
\ee we get
$$
\nabla_iR-2R_{ij}\nabla_jf=0.
$$
Then
$$
\nabla_i(|\nabla f|^2+R)=2\nabla_jf(\nabla_i\nabla_jf+R_{ij})=0.
$$
Therefore \be
R+|\nabla f|^2=C %\eqno(1.1.15)
\ee for some constant $C$.

Taking the difference of (1.1.13) and (1.1.15), we get \be
\Delta f -|\nabla f|^2=-C. %\eqno(1.1.16)
\ee We claim $C=0$ when $M$ is compact. Indeed, this follows
either from \be 0=-\int_M \Delta(e^{-f}) dV
=\int_M (\Delta f-|\nabla f|^2)e^{-f}dV, %\eqno(1.1.17)
\ee or from considering (1.1.16) at both the maximum point and
minimum point of $f$. Then, by integrating (1.1.16) we obtain
$$
\int_M |\nabla f|^2 dV=0.
$$
Therefore $f$ is a constant and $g_{ij}$ is Ricci flat.

For the expanding case, write the soliton equation as
\[
R_{ij}+\nabla_i\nabla_jf-\lambda g_{ij}=0,
\qquad \lambda<0.
\]
Tracing gives $R+\Delta f=n\lambda$.  The differentiated soliton
equation and the contracted Bianchi identity, exactly as above, give
$\nabla_iR=2R_{ij}\nabla_jf$.  It follows that
\[
\nabla_i\bigl(R+|\nabla f|^2-2\lambda f\bigr)=0,
\]
so $R+|\nabla f|^2-2\lambda f=C$ for a constant $C$.  Subtracting
this identity from the trace identity yields
\[
\Delta f-|\nabla f|^2+2\lambda f=n\lambda-C.
\]
At a maximum point of $f$ the right-hand side is at most
$2\lambda\max f$, while at a minimum point it is at least
$2\lambda\min f$.  Hence
\[
2\lambda\min f\le n\lambda-C\le2\lambda\max f.
\]
Since $\lambda<0$ and $\min f\le\max f$, these inequalities force
$\min f=\max f$.  Thus $f$ is constant, and the soliton equation
reduces to $R_{ij}=\lambda g_{ij}$, so the metric is Einstein.
%%%$$\eqno \#$$
\end{proof}


%\vskip 0.3cm \noindent {\bf Remark 1.1.2}
% CZ-BASELINE source-line:850 source-kind:remark source-id:1.1.2
\begin{remark}[examples of nontrivial Ricci solitons]
  \label{rem:cz-examples-of-nontrivial-ricci-solitons}
  \dcref{Remark 1.1.2}
  \group{Cao--Zhu The Ricci Flow}
  \source{cao-zhu}
By contrast, there do exist nontrivial compact gradient shrinking
Ricci solitons (see Koiso \cite{Ko}, Cao \cite{Cao94} and Wang-Zhu
\cite{WZ} ). Also, there exist complete noncompact steady gradient
Ricci solitons that are not Ricci flat. In two dimensions Hamilton
\cite{Ha88} wrote down the first such example on $\mathbb R^2$,
called the {\bf cigar soliton}\index{soliton!cigar}, where the
metric is given by \be
ds^2=\frac {dx^2 +dy^2} {1+x^2+y^2}, %\eqno(1.1.18)
\ee and the vector field is radial, given by $V=-\partial/\partial
r=-(x\partial/\partial x+y\partial/\partial y).$ This metric has
positive curvature and is asymptotic to a cylinder of finite
circumference $2\pi$ at $\infty$. Higher dimensional examples were
found by Robert Bryant \cite{BR} on $\mathbb R^n$ in the
Riemannian case, and by the first author \cite{Cao94} on $\mathbb
C^n$ in the K\"ahler case. These examples are complete,
rotationally symmetric, of positive curvature and found by solving
certain nonlinear ODE (system). Noncompact expanding solitons were
also constructed by the first author \cite{Cao94}. More recently,
Feldman, Ilmanen and Knopf \cite{FIM} constructed new examples of
noncompact shrinking and expanding K\"ahler-Ricci solitons.
\end{remark}


\section{Short-time Existence and Uniqueness}

In this section we establish the short-time existence and
uniqueness result \cite{Ha82, De} for the Ricci flow (1.1.5) on a
compact $n$-dimensional manifold $M$. Our presentation follows
closely Hamilton \cite{Ha82} and DeTurck \cite{De}.

We will see that the Ricci flow is a system of second order
nonlinear weakly parabolic partial differential equations. In
fact, as observed by Hamilton \cite{Ha82}, the degeneracy of the
system is caused by the diffeomorphism group of $M$ which acts as
the gauge group of the Ricci flow. For any diffeomorphism
$\varphi$ of $M$, we have $\Ric(\varphi^*(g))=\varphi^*(\Ric(g))$.
Thus, if $g(t)$ is a solution to the Ricci flow (1.1.5), so is
$\varphi^*(g(t))$. Because the Ricci flow (1.1.5) is only weakly
parabolic, even the existence and uniqueness result on a compact
manifold does not follow from standard PDE theory. The short-time
existence and uniqueness result in the compact case is first
proved by Hamilton \cite{Ha82} using the Nash-Moser implicit
function theorem. Shortly after Denis De Turck \cite{De} gave a
much simpler proof using the gauge fixing idea which we will
present here. In the noncompact case, the short-time existence was
established by Shi \cite{Sh89} in 1989, but the uniqueness result
has been proved only very recently by Bing-Long Chen and the
second author. These results will be presented at the end of this
section.

First, let us follow Hamilton \cite{Ha82} to examine the equation
of the Ricci flow and see it is weakly parabolic. Let $M$ be a
compact $n$-dimensional Riemannian manifold. The Ricci flow
equation is a second order nonlinear partial differential system
\be
\frac{\partial}{\partial t}g_{ij}=E(g_{ij}),%\eqno(1.2.1)
\ee for a family of Riemannian metrics $g_{ij}(\cdot,t)$ on $M$,
where
\begin{align*}
E(g_{ij})&  =  -2R_{ij}\\
&  =  -2\left(\frac{\partial}{\partial x^k}
\Gamma^k_{ij}-\frac{\partial}{\partial x^i}\Gamma^k_{kj}
+\Gamma^k_{kp}\Gamma^p_{ij}-\Gamma^k_{ip}\Gamma^p_{kj}\right)\\
&  = \frac{\partial}{\partial x^i}
\left\{g^{kl}\frac{\partial}{\partial x^j}g_{kl}\right\}
-\frac{\partial}{\partial x^k}
\left\{g^{kl}\left(\frac{\partial}{\partial x^i}g_{jl}
+\frac{\partial}{\partial x^j}g_{il}
-\frac{\partial}{\partial x^l}g_{ij}\right)\right\}\\
&\quad  +2\Gamma^k_{ip}\Gamma^p_{kj}-2\Gamma^k_{kp}\Gamma^p_{ij}.
\end{align*}
The linearization of this system is
$$
\frac{\partial \tilde{g}_{ij}}{\partial t}=DE(g_{ij})
\tilde{g}_{ij}
$$
where $\tilde{g}_{ij}$ is the variation in $g_{ij}$ and $DE$ is
the derivative of $E$ given by
\begin{align*}
DE(g_{ij})\tilde{g}_{ij}&  = g^{kl}\left\{\frac{\partial^2
\tilde{g}_{kl}}{\partial x^i \partial x^j}-\frac{\partial^2
\tilde{g}_{jl}}{\partial x^i \partial x^k}-\frac{\partial^2
\tilde{g}_{il}}{\partial x^j \partial x^k}+\frac{\partial^2
\tilde{g}_{ij}}{\partial x^k \partial x^l}\right\}\\
& \quad + \text{(lower order terms).}
\end{align*}
We now compute the symbol of $DE$. This is to take the highest
order derivatives and replace $\frac{\partial}{\partial x^i}$ by
the Fourier transform variable $\zeta_i$. The symbol of the linear
differential operator $DE(g_{ij})$ in the direction
$\zeta=(\zeta_1 , \ldots,\zeta_n )$ is
$$
\sigma DE(g_{ij})(\zeta)\tilde{g}_{ij} =g^{kl}(\zeta_i \zeta_j
\tilde{g}_{kl} + \zeta_k \zeta_l \tilde{g}_{ij} -\zeta_i \zeta_k
\tilde{g}_{jl} -\zeta_j \zeta_k \tilde{g}_{il}).
$$
To see what the symbol does, we can always assume $\zeta$ has
length 1 and choose coordinates at a point such that
$$
      \left\{
       \begin{array}{lll}
       \displaystyle
g_{ij}=\delta_{ij},
          \\[4mm]
      \displaystyle
\zeta=(1,0,\ldots,0).
       \end{array}
      \right.
$$
Then
\begin{align*}
(\sigma DE(g_{ij})(\zeta))(\tilde{g}_{ij}) &  =
\tilde{g}_{ij}+\delta_{i1}\delta_{j1}(\tilde{g}_{11}
+\cdots+\tilde{g}_{nn})\\
&\quad -\delta_{i1}\tilde{g}_{1j}-\delta_{j1}\tilde{g}_{1i},
\end{align*}
i.e.,
\begin{align*}
[\sigma DE(g_{ij})(\zeta)(\tilde{g}_{ij})]_{11}
&  = \tilde{g}_{22}+\cdots+\tilde{g}_{nn} ,\\
[\sigma DE(g_{ij})(\zeta)(\tilde{g}_{ij})]_{1k} &  = 0 ,
\;\;\;\;\;\; \mbox{ if }\; k\neq1,\\
[\sigma DE(g_{ij})(\zeta)(\tilde{g}_{ij})]_{kl} & =
\tilde{g}_{kl}, \;\;\;\;\mbox{ if }\; k\neq1 ,l\neq1.
\end{align*}
In particular
$$
(\tilde{g}_{ij})=\left(
\begin{array}{cccc}\ast&  \ast&  \cdots&  \ast\\
\ast&  0&  \cdots&  0\\
\vdots&  \vdots&  \ddots&  \vdots\\
\ast&  0&  \cdots&  0
\end{array}
\right)
$$
are zero eigenvectors of the symbol. The presence of the zero
eigenvalue shows that the system cannot be strictly parabolic.

Next, instead of considering the system (1.2.1) (or the Ricci flow
equation (1.1.5)), we follow a trick of DeTurck \cite{De} to
consider a modified evolution equation, which is equivalent to the
Ricci flow up to diffeomorphisms and turns out to be strictly
parabolic so that the standard theory of parabolic equations
applies.

Suppose $\hat{g}_{ij}(x,t)$ is a solution of the Ricci flow
(1.1.5), and $\varphi_t: M\rightarrow M$ is a family of
diffeomorphisms of $M$. Let
$$
g_{ij}(x,t)=\varphi^\ast_t\hat{g}_{ij}(x,t)
$$
be the pull-back metrics. We now want to find the evolution
equation for the metrics $g_{ij}(x,t)$.

Denote by
$$
y(x,t)=\varphi_t(x)=\{y^1(x,t),y^2(x,t),\ldots,y^n(x,t)\}
$$
in local coordinates. Then \be g_{ij}(x,t) =\frac{\partial
y^\alpha}{\partial x^i} \frac{\partial y^\beta}{\partial
x^j}\hat{g}_{\alpha\beta}(y,t),
%\eqno(1.2.2)
\ee and
\begin{align*}
\frac{\partial}{\partial t}g_{ij}(x,t) & =\frac{\partial}{\partial
t}\left[\frac{\partial y^\alpha}{\partial x^i}\frac{\partial
y^\beta}{\partial x^j}
\hat{g}_{\alpha\beta}(y,t)\right]\\
&  = \frac{\partial y^\alpha}{\partial x^i} \frac{\partial
y^\beta}{\partial x^j}\frac{\partial}{\partial t}
\hat{g}_{\alpha\beta}(y,t)+\frac{\partial}{\partial x^i}
\left(\frac{\partial y^\alpha}{\partial t}\right)
\frac{\partial y^\beta}{\partial x^j}\hat{g}_{\alpha\beta}(y,t)\\
&\quad +\frac{\partial y^\alpha}{\partial x^i}
\frac{\partial}{\partial x^j}\left(\frac{\partial
y^\beta}{\partial t}\right) \hat{g}_{\alpha\beta}(y,t).
\end{align*}
Let us choose a normal coordinate $\{x^i\}$ around a fixed point
$p\in M$ such that $\frac{\partial g_{ij}}{\partial x^k}=0$ at
$p$. Since
$$
\frac{\partial}{\partial t}\hat{g}_{\alpha\beta}(y,t)
=-2\hat{R}_{\alpha\beta}(y,t)+\frac{\partial
\hat{g}_{\alpha\beta}}{\partial y^\gamma}\frac{\partial
y^\gamma}{\partial t} ,
$$
we have in the normal coordinate,
\begin{align*}
&\frac{\partial}{\partial t}g_{ij}(x,t) \\
&  =-2\frac{\partial y^\alpha}{\partial x^i}\frac{\partial
y^\beta}{\partial x^j}\hat{R}_{\alpha\beta}(y,t)+\frac{\partial
y^\alpha}{\partial x^i}\frac{\partial y^\beta}{\partial
x^j}\frac{\partial \hat{g}_{\alpha\beta}}{\partial
y^\gamma}\frac{\partial y^\gamma}{\partial t}\\
&\quad +\frac{\partial}{\partial x^i}\left(\frac{\partial
y^\alpha}{\partial t}\right)\frac{\partial y^\beta}{\partial
x^j}\hat{g}_{\alpha\beta}(y,t)+\frac{\partial}{\partial
x^j}\left(\frac{\partial y^\beta}{\partial t}\right)\frac{\partial
y^\alpha}{\partial x^i}\hat{g}_{\alpha\beta}(y,t)\\
&  = -2R_{ij}(x,t)+\frac{\partial y^\alpha}{\partial
x^i}\frac{\partial y^\beta}{\partial x^j}\frac{\partial
\hat{g}_{\alpha\beta}}{\partial y^\gamma}\frac{\partial
y^\gamma}{\partial t}+\frac{\partial}{\partial
x^i}\left(\frac{\partial y^\alpha}{\partial
t}\right)\frac{\partial x^k}{\partial
y^\alpha}g_{jk} \\
&\quad+\frac{\partial}{\partial x^j}\left(\frac{\partial
y^\beta}{\partial t}\right)\frac{\partial x^k}{\partial
y^\beta}g_{ik}\\
&  = -2R_{ij}(x,t)+\frac{\partial y^\alpha}{\partial
x^i}\frac{\partial y^\beta}{\partial x^j}\frac{\partial
\hat{g}_{\alpha\beta}}{\partial y^\gamma}\frac{\partial
y^\gamma}{\partial t}+\frac{\partial}{\partial
x^i}\left(\frac{\partial y^\alpha}{\partial t}\frac{\partial
x^k}{\partial
y^\alpha}g_{jk}\right) \\
&\quad+\frac{\partial}{\partial x^j}\left(\frac{\partial
y^\beta}{\partial t}\frac{\partial x^k}{\partial
y^\beta}g_{ik}\right)%\\
%&\quad
-\frac{\partial y^\alpha}{\partial t}\frac{\partial}{\partial
x^i}\left(\frac{\partial x^k}{\partial
y^\alpha}\right)g_{jk}-\frac{\partial y^\beta}{\partial
t}\frac{\partial}{\partial x^j}\left(\frac{\partial x^k}{\partial
y^\beta}\right)g_{ik}.
\end{align*}
The second term on the RHS gives, in the normal coordinate,
\begin{align*}
\frac{\partial y^\alpha}{\partial x^i}\frac{\partial
y^\beta}{\partial x^j}\frac{\partial y^\gamma}{\partial
t}\frac{\partial \hat{g}_{\alpha\beta}}{\partial y^\gamma} &
=\frac{\partial y^\alpha}{\partial x^i}\frac{\partial
y^\beta}{\partial x^j}\frac{\partial y^\gamma}{\partial
t}g_{kl}\frac{\partial}{\partial y^\gamma} \(\frac{\partial
x^k}{\partial y^\alpha}
\frac{\partial x^l}{\partial y^\beta}\)\\
&  = \frac{\partial y^\alpha}{\partial x^i}\frac{\partial
y^\gamma}{\partial t}\frac{\partial}{\partial
y^\gamma}\(\frac{\partial x^k}{\partial
y^\alpha}\)g_{jk}+\frac{\partial y^\beta}{\partial
x^j}\frac{\partial y^\gamma}{\partial t}\frac{\partial}{\partial
y^\gamma}\(\frac{\partial x^k}{\partial y^\beta}\)g_{ik}\\
&  = \frac{\partial y^\alpha}{\partial t}\frac{\partial^2
x^k}{\partial y^\alpha\partial y^\beta}\frac{\partial
y^\beta}{\partial x^i}g_{jk}+\frac{\partial y^\beta}{\partial
t}\frac{\partial^2 x^k}{\partial y^\alpha\partial
y^\beta}\frac{\partial y^\alpha}{\partial x^j}g_{ik}\\
&  = \frac{\partial y^\alpha}{\partial t}\frac{\partial}{\partial
x^i}\(\frac{\partial x^k}{\partial
y^\alpha}\)g_{jk}+\frac{\partial y^\beta}{\partial
t}\frac{\partial}{\partial x^j}\(\frac{\partial x^k}{\partial
y^\beta}\)g_{ik}.
\end{align*}
So we get
\begin{align}
&\frac{\partial}{\partial t}g_{ij}(x,t) \\
&=-2R_{ij}(x,t)+\nabla_i\left(\frac{\partial y^\alpha}{\partial
t}\frac{\partial x^k}{\partial
y^\alpha}g_{jk}\right)+\nabla_j\left(\frac{\partial
y^\beta}{\partial
t}\frac{\partial x^k}{\partial y^\beta}g_{ik}\right).\nn%\eqno(1.2.3)
\end{align}
If we define $y(x,t)=\varphi_t(x)$ by the equations \be
      \left\{
       \begin{array}{lll}
  \frac{\partial y^\alpha}{\partial t}=\frac{\partial y^\alpha}{\partial
x^k}g^{jl}(\Gamma^k_{jl}-\stackrel{o}{\Gamma}^k_{jl}),
          \\[4mm]
  y^\alpha(x,0)=x^\alpha,
       \end{array}
    \right. %\eqno(1.2.4)
\ee and
$V_i=g_{ik}g^{jl}(\Gamma^k_{jl}-\stackrel{o}{\Gamma}^k_{jl})$, we
get the following evolution equation for the pull-back metric \be
      \left\{
       \begin{array}{lll}
  \frac{\partial}{\partial t}g_{ij}(x,t)
=-2R_{ij}(x,t)+\nabla_iV_j+\nabla_jV_i, \\[4mm]
  g_{ij}(x,0)= \stackrel{o}{g}_{ij}(x),
       \end{array}
    \right. % \eqno(1.2.5)
\ee where $\stackrel{o}{g}_{ij}(x)$ is the initial metric and
$\stackrel{o}{\Gamma}^k_{jl}$ is the connection of the initial
metric.

%%%\vskip 0.2cm \noindent {\bf Lemma 1.2.1} \emph{
% CZ-BASELINE source-line:1137 source-kind:lemma source-id:1.2.1
\begin{lemma}[strict parabolicity of the DeTurck-modified flow]
  \label{lem:cz-strict-parabolicity-of-the-deturck-modified-flow}
  \dcref{Lemma 1.2.1}
  \group{Cao--Zhu Short-time Existence and Uniqueness}
  \source{cao-zhu}
The modified evolution equation $(1.2.5)$ is a strictly parabolic
system.
\end{lemma}

%%%\vskip 0.1cm \noindent {\bf Proof.}
\begin{proof}
The RHS of the equation (1.2.5) is given by
\begin{align*}
&  -2R_{ij}(x,t)+\nabla_iV_j+\nabla_jV_i\\
&= \frac{\partial}{\partial x^i}\left\{g^{kl}\frac{\partial
g_{kl}}{\partial x^j}\right\}-\frac{\partial}{\partial
x^k}\left\{g^{kl}\(\frac{\partial g_{jl}}{\partial
x^i}+\frac{\partial
g_{il}}{\partial x^j}-\frac{\partial g_{ij}}{\partial x^l}\)\right\}\\
&\quad  +g_{jk}g^{pq}\frac{\partial}{\partial
x^i}\left\{\frac{1}{2}g^{kl}\(\frac{\partial g_{pl}}{\partial
x^q}+\frac{\partial g_{ql}}{\partial
x^p}-\frac{\partial g_{pq}}{\partial x^l}\)\right\}\\
&\quad  +g_{ik}g^{pq}\frac{\partial}{\partial
x^j}\left\{\frac{1}{2}g^{kl}\(\frac{\partial g_{pl}}{\partial
x^q}+\frac{\partial g_{ql}}{\partial
x^p}-\frac{\partial g_{pq}}{\partial x^l}\)\right\}\\
&\quad  +\text{(lower order terms)}\\
&= g^{kl}\left\{\frac{\partial^2 g_{kl}}{\partial x^i\partial
x^j}-\frac{\partial^2 g_{jl}}{\partial x^i\partial
x^k}-\frac{\partial^2 g_{il}}{\partial x^j\partial
x^k}+\frac{\partial^2 g_{ij}}{\partial x^k\partial
x^l}\right\}\\
&\quad  +\frac{1}{2}g^{pq}\left\{\frac{\partial^2 g_{pj}}{\partial
x^i\partial x^q}+\frac{\partial^2 g_{qj}}{\partial x^i\partial
x^p}-\frac{\partial^2
g_{pq}}{\partial x^i\partial x^j}\right\}\\
&\quad  +\frac{1}{2}g^{pq}\left\{\frac{\partial^2 g_{p
i}}{\partial x^j\partial x^q}+\frac{\partial^2 g_{q i}}{\partial
x^j\partial x^p}-\frac{\partial^2
g_{pq}}{\partial x^i\partial x^j}\right\}\\
&\quad +\text{(lower order terms)}\\
&= g^{kl}\frac{\partial^2 g_{ij}}{\partial x^k\partial x^l}
+\text{(lower order terms)}.
\end{align*}
Thus its symbol is $(g^{kl}\zeta_k\zeta_l)\tilde{g}_{ij}$. Hence
the equation in (1.2.5) is strictly parabolic.
%%$$\eqno \#$$ \vskip 0.3cm
\end{proof}

Now since the equation (1.2.5) is strictly parabolic and the
manifold $M$ is compact,  it follows from the standard theory of
parabolic equations (see for example \cite{LSU}) that (1.2.5) has
a solution for a short time.  From the solution of (1.2.5) we can
obtain a solution of the Ricci flow from (1.2.4) and (1.2.2). This
shows existence. Now we argue the uniqueness of the solution.
Since
$$
\Gamma^k_{jl} = \frac{\partial y^{\alpha}}{\partial x^j}
\frac{\partial y^{\beta}}{\partial x^l} \frac{\partial
x^k}{\partial y^{\gamma}}\hat{\Gamma}^{\gamma}_{\alpha \beta} +
\frac{\partial x^k}{\partial y^{\alpha}} \frac{\partial ^2
y^{\alpha}}{\partial x^j \partial x^l},
$$
the initial value problem (1.2.4) can be written as \be
\begin{cases}
\frac{\partial y^\alpha}{\partial t} =g^{jl}\(\frac{\partial ^2
y^\alpha}{\partial x^j \partial
x^l}-\stackrel{o}{\Gamma}^k_{jl}\frac{\partial
y^{\alpha}}{\partial x^k}+\hat{\Gamma}^{\alpha}_{\gamma
\beta}\frac{\partial y^{\beta}}{\partial x^j}\frac{\partial
y^{\gamma}}{\partial x^l}\),\\[4mm]
  y^\alpha(x,0)=x^\alpha.
       \end{cases} %\eqno(1.2.6)
\ee This is clearly a strictly parabolic system. For any two
solutions $\hat{g}^{(1)}_{ij}(\cdot,t)$ and
$\hat{g}^{(2)}_{ij}(\cdot,t)$ of the Ricci flow (1.1.5) with the
same initial data, we can solve the initial value problem (1.2.6)
(or equivalently, (1.2.4)) to get two families $\varphi ^{(1)}_t$
and $\varphi ^{(2)}_t$ of diffeomorphisms of $M$. Thus we get two
solutions, ${g}^{(1)}_{ij}(\cdot,t) = (\varphi
^{(1)}_t)^*\hat{g}^{(1)}_{ij}(\cdot,t)$ and
${g}^{(2)}_{ij}(\cdot,t) = (\varphi
^{(2)}_t)^*\hat{g}^{(2)}_{ij}(\cdot,t)$, to the modified evolution
equation (1.2.5) with the same initial metric. The uniqueness
result for the strictly parabolic equation (1.2.5) implies that
$g^{(1)}_{ij} = g^{(2)}_{ij}$. Then by equation (1.2.4) and the
standard uniqueness result of ODE systems, the corresponding
solutions $\varphi ^{(1)}_t$ and $\varphi ^{(2)}_t$ of (1.2.4) (or
equivalently (1.2.6)) must agree. Consequently the metrics
$\hat{g}^{(1)}_{ij}$ and $\hat{g}^{(2)}_{ij}$ must agree also.
Thus we have proved the following result.

%%\vskip 0.2cm \noindent {\bf Theorem 1.2.2} (Hamilton \cite{Ha82},
%%De Turck \cite{De}) \emph{
% CZ-BASELINE source-line:1228 source-kind:theorem source-id:1.2.2
\begin{theorem}[short-time existence and uniqueness on compact manifolds]
  \label{thm:cz-short-time-existence-and-uniqueness-on-compact-manifolds}
  \dcref{Theorem 1.2.2}
  \group{Cao--Zhu Short-time Existence and Uniqueness}
  \source{cao-zhu}
Let $(M,\, g_{ij}(x))$ be a compact Riemannian manifold. Then
there exists a constant $T>0$ such that the initial value problem
$$
\begin{cases}
  \frac{\partial}{\partial t}g_{ij}(x,t)=-2R_{ij}(x,t)
          \\[4mm]
  g_{ij}(x,0)=g_{ij}(x)
       \end{cases}
$$
has a unique smooth solution $g_{ij}(x,t)$ on $M\times [0,T)$.
%%$$ \eqno \#$$ \vskip 0.3cm
\end{theorem}

The case of a noncompact manifold is much more complicated and
involves a huge amount of techniques from the theory of partial
differential equations. Here we will only state the existence and
uniqueness results and refer the reader to the cited references
for the proofs.

The following existence result was obtained by Shi in \cite{Sh89}
published in 1989.

%%\vskip 0.2cm \noindent {\bf Theorem 1.2.3} (Shi
%%\cite{Sh89}) \emph{
% CZ-BASELINE source-line:1253 source-kind:theorem source-id:1.2.3
\begin{theorem}[short-time existence on complete noncompact manifolds]
  \label{thm:cz-short-time-existence-on-complete-noncompact-manifolds}
  \dcref{Theorem 1.2.3}
  \group{Cao--Zhu Short-time Existence and Uniqueness}
  \source{cao-zhu}
Let $(M,\, g_{ij}(x))$ be a complete noncompact Riemannian
manifold of dimension $n$ with bounded curvature. Then there
exists a constant $T>0$ such that the initial value problem
$$
\begin{cases}
  \frac{\partial}{\partial t}g_{ij}(x,t)=-2R_{ij}(x,t)
          \\[4mm]
  g_{ij}(x,0)=g_{ij}(x)
       \end{cases}
$$
has a smooth solution $g_{ij}(x,t)$ on $M\times [0,T]$ with
uniformly bounded curvature.
%%$$ \eqno \#$$ \vskip 0.3cm
\end{theorem}

The Ricci flow is a heat type equation. It is well-known that the
uniqueness of a heat equation on a complete noncompact manifold is
not always held if there are no further restrictions on the growth
of the solutions. For example, the heat equation on Euclidean
space with zero initial data has a nontrivial solution which grows
faster than $\exp(a|x|^2)$ for any $a>0$ whenever $t>0$. This
implies that even for the standard linear heat equation on
Euclidean space, in order to ensure the uniqueness one can only
allow the solution to grow at most as $\exp(C|x|^2)$ for some
constant $C>0$. Note that on a K\"ahler manifold, the Ricci
curvature is given by $R_{\alpha\bar\beta} = - \frac{\partial
^2}{\partial z^{\alpha}\partial \bar{z}^{\beta}}
\log\det(g_{\gamma\bar{\delta}})$. So the reasonable growth rate
for the uniqueness of the Ricci flow to hold is that the solution
has bounded curvature. Thus the following uniqueness result of
Bing-Long Chen and the second author \cite{CZ05U} is essentially
the best one can hope for.

%%\vskip 0.2cm\noindent {\bf Theorem 1.2.4} (Chen-Zhu \cite{CZ05U})
%%\emph{
% CZ-BASELINE source-line:1289 source-kind:theorem source-id:1.2.4
\begin{theorem}[bounded-curvature uniqueness on complete noncompact manifolds]
  \label{thm:cz-bounded-curvature-uniqueness-on-complete-noncompact-manifolds}
  \dcref{Theorem 1.2.4}
  \group{Cao--Zhu Short-time Existence and Uniqueness}
  \source{cao-zhu}
Let $(M, \hat{g}_{ij})$ be a complete noncompact Riemannian
manifold of dimension $n$ with bounded curvature. Let
${g}_{ij}(x,t)$ and $\bar{g}_{ij}(x,t)$ be two solutions, defined
on $M\times[0,T]$, to the Ricci flow $(1.1.5)$ with $\hat{g}_{ij}$
as initial data and with bounded curvatures. Then
$g_{ij}(x,t)\equiv \bar{g}_{ij}(x,t)$ on $M\times[0,T]$.
%%$$\eqno \#$$ \vskip 1cm
\end{theorem}


\section{Evolution of Curvatures}

The Ricci flow is an evolution equation on the metric. The
evolution for the metric implies a nonlinear heat equation for the
Riemannian curvature tensor $R_{ijkl}$ which we now derive. Our
presentation in this section follows closely the original papers
of Hamilton \cite{Ha82, Ha86}.

%%\vskip 0.2cm\noindent{\bf Proposition 1.3.1} (Hamilton
%%\cite{Ha82}) \emph{
% CZ-BASELINE source-line:1310 source-kind:proposition source-id:1.3.1
\begin{proposition}[evolution of the Riemann curvature tensor]
  \label{prop:cz-evolution-of-the-riemann-curvature-tensor}
  \dcref{Proposition 1.3.1}
  \group{Cao--Zhu Evolution of Curvatures}
  \source{cao-zhu}
Under the Ricci flow $(1.1.5),$ the curvature tensor satisfies the
evolution equation
\begin{align*}
\frac{\partial}{\partial t}R_{ijkl}&  = \Delta
R_{ijkl}+2(B_{ijkl}-B_{ijlk}-B_{iljk}+B_{ikjl})\\
& \quad
-g^{pq}(R_{pjkl}R_{qi}+R_{ipkl}R_{qj}+R_{ijpl}R_{qk}+R_{ijkp}R_{ql})
\end{align*}
where $B_{ijkl}=g^{pr}g^{qs}R_{piqj}R_{rksl}$ and $\Delta$ is the
Laplacian with respect to the evolving metric.
\end{proposition}

%%\vskip 0.1cm\noindent {\bf Proof.}\
\begin{proof}
Choose $\{x^1,\ldots,x^m\}$ to be a normal coordinate system at a
fixed point. At this point, we compute
\begin{align*}
\frac{\partial}{\partial t}\Gamma^h_{jl} &
=\frac{1}{2}g^{hm}\left\{\frac{\partial}{\partial
x^j}\(\frac{\partial}{\partial t}g_{lm}\)+\frac{\partial}{\partial
x^l}\(\frac{\partial}{\partial t}g_{jm}\)-\frac{\partial}{\partial
x^m}\(\frac{\partial}{\partial t}g_{jl}\)\right\}\\
&  = \frac{1}{2}g^{hm}(\nabla_j(-2R_{lm})
+\nabla_l(-2R_{jm})-\nabla_m(-2R_{jl})),\\
\frac{\partial}{\partial t}R_{ijl}^h & = \frac{\partial}{\partial
x^i}\(\frac{\partial}{\partial
t}\Gamma^h_{jl}\)-\frac{\partial}{\partial
x^j}\(\frac{\partial}{\partial t}\Gamma^h_{il}\),\\
\frac{\partial}{\partial t}R_{ijkl} &
=g_{hk}\frac{\partial}{\partial t}R^h_{ijl}+\frac{\partial
g_{hk}}{\partial t}R^h_{ijl}.
\end{align*}

\noindent Combining these identities we get
\begin{align*}
\frac{\partial}{\partial t}R_{ijkl} &
=g_{hk}\bigg[\(\frac{1}{2}\nabla_i
[g^{hm}(\nabla_j(-2R_{lm})+\nabla_l(-2R_{jm})-\nabla_m(-2R_{jl}))]\)\\
&\quad -\(\frac{1}{2}\nabla_j[g^{hm}(\nabla_i(-2R_{lm})
+\nabla_l(-2R_{im})-\nabla_m(-2R_{il}))]\)\bigg]\\
& \quad -2R_{hk}R^h_{ijl}\\
&  = \nabla_i\nabla_kR_{jl}-\nabla_i\nabla_lR_{jk}
-\nabla_j\nabla_kR_{il}+\nabla_j\nabla_lR_{ik}\\
&\quad -R_{ijlp}g^{pq}R_{qk}-R_{ijkp}g^{pq}R_{ql}-2R_{ijpl}g^{pq}R_{qk}\\
& = \nabla_i\nabla_kR_{jl}-\nabla_i\nabla_lR_{jk}
-\nabla_j\nabla_kR_{il}+\nabla_j\nabla_lR_{ik}\\
&\quad -g^{pq}(R_{ijkp}R_{ql}+R_{ijpl}R_{qk}).
\end{align*}
Here we have used the exchanging formula (1.1.3).

Now it remains to check the following identity, which is analogous
to the Simon$'$s identity in extrinsic geometry,
\begin{align}
& \Delta R_{ijkl}+2(B_{ijkl}-B_{ijlk}-B_{iljk}+B_{ikjl})\\
&  = \nabla_i\nabla_k R_{jl}-\nabla_i\nabla_l
R_{jk}-\nabla_j\nabla_k R_{il}+\nabla_j\nabla_l R_{ik} \nn\\
&\quad +g^{pq}(R_{pjkl}R_{qi}+R_{ipkl}R_{qj}). \nn
%\eqno(1.3.1)
\end{align}
Indeed, from the second Bianchi identity (1.1.4), we have
\begin{align*}
\Delta R_{ijkl}&  = g^{pq}\nabla_p\nabla_q R_{ijkl}\\
&  = g^{pq}\nabla_p\nabla_i
R_{qjkl}-g^{pq}\nabla_p\nabla_jR_{qikl}.
\end{align*}
Let us examine the first term on the RHS. By using the exchanging
formula (1.1.3) and the first Bianchi identity (1.1.1), we have
\begin{align*}
& g^{pq}\nabla_p\nabla_iR_{qjkl}-g^{pq}\nabla_i\nabla_pR_{qjkl}\\
& = g^{pq}g^{mn}(R_{piqm}R_{njkl}+R_{pijm}R_{qnkl}
+R_{pikm}R_{qjnl}+R_{pilm}R_{qjkn})\\
& = R_{im}g^{mn}R_{njkl}+g^{pq}g^{mn}R_{pimj}(R_{qkln}+R_{qlnk})\\
&\quad +g^{pq}g^{mn}R_{pikm}R_{qjnl}+g^{pq}g^{mn}R_{pilm}R_{qjkn}\\
& = R_{im}g^{mn}R_{njkl}-B_{ijkl}+B_{ijlk}-B_{ikjl}+B_{iljk},
\end{align*}
while using the contracted second Bianchi identity \be
g^{pq}\nabla_pR_{qjkl}=\nabla_kR_{jl}-\nabla_lR_{jk},%\eqno(1.3.2)
\ee we have
$$
g^{pq}\nabla_i\nabla_pR_{qjkl}
=\nabla_i\nabla_kR_{jl}-\nabla_i\nabla_lR_{jk}.
$$
Thus
\begin{align*}
& g^{pq}\nabla_p\nabla_iR_{qjkl} \\
&
=\nabla_i\nabla_kR_{jl}-\nabla_i\nabla_lR_{jk}-(B_{ijkl}-B_{ijlk}
-B_{iljk}+B_{ikjl})+g^{pq}R_{pjkl}R_{qi}.
\end{align*}
Therefore we obtain
\begin{align*}
& \Delta R_{ijkl} \\
&  = g^{pq}\nabla_p\nabla_iR_{qjkl}-g^{pq}\nabla_p\nabla_jR_{qikl}\\
&
=\nabla_i\nabla_kR_{jl}-\nabla_i\nabla_lR_{jk}-(B_{ijkl}-B_{ijlk}
-B_{iljk}+B_{ikjl})+g^{pq}R_{pjkl}R_{qi}\\
&\quad
-\nabla_j\nabla_kR_{il}+\nabla_j\nabla_lR_{ik}+(B_{jikl}-B_{jilk}
-B_{jlik}+B_{jkil})-g^{pq}R_{pikl}R_{qj}\\
& = \nabla_i\nabla_kR_{jl}-\nabla_i\nabla_lR_{jk}
-\nabla_j\nabla_kR_{il}+\nabla_j\nabla_lR_{ik}\\[2mm]
&\quad +g^{pq}(R_{pjkl}R_{qi}+R_{ipkl}R_{qj})-2(B_{ijkl}-B_{ijlk}
-B_{iljk}+B_{ikjl})
\end{align*}
as desired, where in the last step we used the symmetries \be
B_{ijkl}=B_{klij}=B_{jilk}.%\eqno(1.3.3)
\ee
%%$$\eqno \#$$ \vskip 0.3cm
\end{proof}

%%\noindent {\bf Corollary 1.3.2}\ \emph{
% CZ-BASELINE source-line:1422 source-kind:corollary source-id:1.3.2
\begin{corollary}[evolution of the Ricci tensor]
  \label{cor:cz-evolution-of-the-ricci-tensor}
  \dcref{Corollary 1.3.2}
  \group{Cao--Zhu Evolution of Curvatures}
  \source{cao-zhu}
The Ricci curvature satisfies the evolution equation
$$
\frac{\partial}{\partial t}R_{ik} =\Delta
R_{ik}+2g^{pr}g^{qs}R_{piqk}R_{rs}-2g^{pq}R_{pi}R_{qk}.
$$
\end{corollary}

%%\vskip 0.1cm\noindent{\bf Proof.}\
\begin{proof}
\begin{align*}
\frac{\partial}{\partial t}R_{ik} &
=g^{jl}\frac{\partial}{\partial t}R_{ijkl}
+\(\frac{\partial}{\partial t}g^{jl}\)R_{ijkl}\\
& = g^{jl}[\Delta R_{ijkl}+2(B_{ijkl}-B_{ijlk}
-B_{iljk}+B_{ikjl})\\
&\quad -g^{pq}(R_{pjkl}R_{qi}+R_{ipkl}R_{qj}
+R_{ijpl}R_{qk}+R_{ijkp}R_{ql})]\\
&\quad -g^{jp}\(\frac{\partial}{\partial t}g_{pq}\)g^{ql}R_{ijkl}\\
& = \Delta
R_{ik}+2g^{jl}(B_{ijkl}-2B_{ijlk})+2g^{pr}g^{qs}R_{piqk}R_{rs}\\
&\quad -2g^{pq}R_{pk}R_{qi}.
\end{align*}

We claim that $g^{jl}(B_{ijkl}-2B_{ijlk})=0$. Indeed by using the
first Bianchi identity, we have
\begin{align*}
g^{jl}B_{ijkl}&  = g^{jl}g^{pr}g^{qs}R_{piqj}R_{rksl}\\
&  = g^{jl}g^{pr}g^{qs}R_{pqij}R_{rskl}\\
&  = g^{jl}g^{pr}g^{qs}(R_{piqj}-R_{pjqi})(R_{rksl}-R_{rlsk})\\
&  = 2g^{jl}(B_{ijkl}-B_{ijlk})
\end{align*}
as desired.

Thus we obtain
$$
\frac{\partial}{\partial t}R_{ik} =\Delta
R_{ik}+2g^{pr}g^{qs}R_{piqk}R_{rs}-2g^{pq}R_{pi}R_{qk}.
$$
%%$$\eqno \#$$ \vskip 0.3cm
\end{proof}

%%\noindent {\bf Corollary 1.3.3}\ \emph{
% CZ-BASELINE source-line:1465 source-kind:corollary source-id:1.3.3
\begin{corollary}[evolution of scalar curvature]
  \label{cor:cz-evolution-of-scalar-curvature}
  \dcref{Corollary 1.3.3}
  \group{Cao--Zhu Evolution of Curvatures}
  \source{cao-zhu}
The scalar curvature satisfies the evolution equation
$$
\frac{\partial R}{\partial t}=\Delta R+2|\Ric|^2.
$$
\end{corollary}

%%\vskip 0.1cm\noindent {\bf Proof.}\
\begin{proof}
\begin{align*}
\frac{\partial R}{\partial t}&  = g^{ik}\frac{\partial
R_{ik}}{\partial t}+\(-g^{ip}\frac{\partial g_{pq}}{\partial
t}g^{qk}\)R_{ik}\\
&  = g^{ik}(\Delta R_{ik}+2g^{pr}g^{qs}R_{piqk}R_{rs}
-2g^{pq}R_{pi}R_{qk})+2R_{pq}R_{ik}g^{ip}g^{qk}\\
&  = \Delta R+2|\Ric|^2.
\end{align*}
%%$$\eqno \#$$ \vskip 0.3cm
\end{proof}

To simplify the evolution equations of curvatures, we will follow
Hamilton \cite{Ha86} and represent the curvature tensors in an
orthonormal frame and evolve the frame so that it remains
orthonormal. More precisely, let us pick an abstract vector bundle
$V$ over $M$ isomorphic to the tangent bundle $TM$. Locally, the
frame $F=\{ F_1,\ldots,F_a,\ldots,F_n\}$ of $V$ is given by
$F_a=F^i_a\frac{\partial}{\partial x^i}$ with the isomorphism
$\{F^i_a\}$. Choose $\{F^i_a\}$ at $t=0$ such that
$F=\{F_1,\ldots,F_a,\ldots,F_n\}$ is an orthonormal frame at
$t=0$, and evolve $\{F_a^i\}$ by the equation
$$
\frac{\partial}{\partial t}F^i_a=g^{ij}R_{jk}F^k_a.
$$
Then the frame $F=\{F_1,\ldots,F_a,\ldots,F_n\}$ will remain
orthonormal for all times since the pull back metric on $V$
$$
h_{ab}=g_{ij}F^i_aF^j_b
$$
remains constant in time. In the following we will use indices
$a,b,\ldots$ on a tensor to denote its components in the evolving
orthonormal frame. In this frame we have the following:
\begin{align*}
R_{abcd}&  = F^i_aF^j_bF^k_cF^l_dR_{ijkl},\\
\Gamma^a_{jb}&  = F^a_i\frac{\partial F_b^i}{\partial
x^j}+\Gamma_{jk}^iF_i^aF_b^k,\quad ((F^a_i)=(F^i_a)^{-1})\\
\nabla_iV^a&  = \frac{\partial}{\partial
x^i}V^a+\Gamma_{ib}^aV^b,\\
\nabla_bV^a&  = F^i_b\nabla_iV^a,
\end{align*}
where $\Gamma_{jb}^a$ is the metric connection of the vector
bundle $V$ with the metric $h_{ab}$. Indeed, by direct
computations,
\begin{align*}
\nabla_iF_b^j&  = \frac{\partial F_b^j}{\partial
x^i}+F_b^k\Gamma_{ik}^j-F_c^j\Gamma_{ib}^c\\
&  = \frac{\partial F_b^j}{\partial x^i}+F_b^k\Gamma_{ik}^j
-F^j_c\(F_k^c\frac{\partial F_b^k}{\partial
x^i}+\Gamma_{ik}^lF^c_lF_b^k\)\\
&  = 0,\\
\nabla_ih_{ab}&  = \nabla_i(g_{ij}F^i_aF^j_b)=0.
\end{align*}
So
$$
\nabla_aV_b=F^i_aF^j_b\nabla_iV^j,\hskip 2.5cm$$ and
\begin{align*}
\Delta R_{abcd}&  = \nabla_l\nabla_lR_{abcd}\\
&  = g^{ij}\nabla_i\nabla_jR_{abcd}\\
&  = g^{ij}F_a^kF_b^lF_c^mF_d^n\nabla_i\nabla_jR_{klmn}.
\end{align*}

In an orthonormal frame $F=\{F_1,\ldots,F_a,\ldots,F_n\}$, the
evolution equations of curvature tensors become
\begin{align}
\frac{\partial}{\partial t}R_{abcd}&=\Delta
R_{abcd}+2(B_{abcd}-B_{abdc}-B_{adbc}+B_{acbd}) \\  %%\eqno(1.3.4)
\frac{\partial}{\partial t}R_{ab}&=\Delta
R_{ab}+2R_{acbd}R_{cd}   \\%%\eqno(1.3.5)
\frac{\partial}{\partial t}R&=\Delta R+2|\Ric|^2  %%\eqno(1.3.6)
\end{align}
where $B_{abcd}=R_{aebf}R_{cedf}$.

Equation (1.3.4) is a reaction-diffusion equation. We can
understand the quadratic terms of this equation better if we think
of the curvature tensor $R_{abcd}$ as a symmetric bilinear form on
the two-forms $\Lambda^2(V)$ given by the formula
$$
Rm(\varphi,\psi)=R_{abcd}\varphi_{ab}\psi_{cd},\qquad \text{for
}\; \varphi,\psi\in\Lambda^2(V).
$$

A two-form $\varphi\in\Lambda^2(V)$ can be regarded as an element
of the Lie algebra $so(n)$ (i.e. the skew-symmetric matrix
$(\varphi _{ab})_{n\times n}$), where the metric on $\Lambda^2(V)$
is given by
$$
\<\varphi,\psi\>=\varphi_{ab}\psi_{ab}
$$
and the Lie bracket is given by
$$
[\varphi,\psi]_{ab}=\varphi_{ac}\psi_{bc}-\psi_{ac}\varphi_{bc}.
$$

Choose an orthonormal basis of $\Lambda^2(V)$
$$
\Phi=\{\varphi^1,\ldots,\varphi^\alpha,\ldots,\varphi^{\frac{n(n-1)}{2}}\}
$$
where $\varphi^\alpha=\{\varphi_{ab}^\alpha\}$. The Lie bracket is
given by
$$
[\varphi^\alpha,\varphi^\beta]=C^{\alpha\beta}_\gamma\varphi^\gamma,
$$
where
$C^{\alpha\beta\gamma}=C^{\alpha\beta}_\sigma\delta^{\sigma\gamma}
=\<[\varphi^\alpha,\varphi^\beta],\varphi^\gamma\>$ are the Lie
structure constants.

Write
$R_{abcd}=M_{\alpha\beta}\varphi_{ab}^\alpha\varphi^\beta_{cd}.$
We now claim that the first part of the quadratic terms in (1.3.4)
is given by \be
2(B_{abcd}-B_{abdc})=M_{\alpha\gamma}M_{\beta\gamma}
\varphi_{ab}^\alpha\varphi^\beta_{cd}.  %%\eqno(1.3.7)
\ee Indeed, by the first Bianchi identity,
\begin{align*}
B_{abcd}-B_{abdc}&  = R_{aebf}R_{cedf}-R_{aebf}R_{decf}\\
&  = R_{aebf}(-R_{cefd}-R_{cfde})\\
&  = R_{aebf}R_{cdef}.
\end{align*}
On the other hand,
\begin{align*}
R_{aebf}R_{cdef}&  = (-R_{abfe}-R_{afeb})R_{cdef}\\
&  = R_{abef}R_{cdef}-R_{afeb}R_{cdef}\\
&  = R_{abef}R_{cdef}-R_{afbe}R_{cdfe}
\end{align*}
which implies $R_{aebf}R_{cdef}=\frac{1}{2}R_{abef}R_{cdef}$. Thus
we obtain
$$
2(B_{abcd}-B_{abdc}) =R_{abef}R_{cdef}
=M_{\alpha\gamma}M_{\beta\gamma}\varphi_{ab}^\alpha\varphi^\beta_{cd}.
$$

We next consider the last part of the quadratic terms:
\begin{align*}
& 2(B_{acbd}-B_{adbc}) \\
&  = 2(R_{aecf}R_{bedf}-R_{aedf}R_{becf})\\
&  = 2(M_{\gamma\delta}\varphi_{ae}^\gamma\varphi_{cf}^\delta
M_{\eta\theta}\varphi_{be}^\eta\varphi^\theta_{df}
-M_{\gamma\theta}\varphi_{ae}^\gamma\varphi_{df}^\theta
M_{\eta\delta}\varphi_{be}^\eta\varphi_{cf}^\delta)\\
& =2[M_{\gamma\delta}(\varphi_{ae}^\eta\varphi_{be}^\gamma
+C_\alpha^{\gamma\eta}\varphi_{ab}^\alpha)\varphi_{cf}^\delta
M_{\eta\theta}\varphi^\theta_{df}
-M_{\eta\theta}\varphi^\eta_{ae}\varphi^\theta_{df}
M_{\gamma\delta}\varphi^\gamma_{be}\varphi^\delta_{cf}]\\
&  = 2M_{\gamma\delta}\varphi_{cf}^\delta
M_{\eta\theta}\varphi^\theta_{df}C_\alpha^{\gamma\eta}\varphi_{ab}^\alpha.
\end{align*}
But
\begin{align*}
& M_{\gamma\delta}\varphi_{cf}^\delta
M_{\eta\theta}\varphi^\theta_{df}C_\alpha^{\gamma\eta}
\varphi_{ab}^\alpha\\
&
=M_{\gamma\delta}M_{\eta\theta}C_\alpha^{\gamma\eta}\varphi_{ab}^\alpha
[\varphi^\theta_{cf}\varphi^\delta_{df}
+C_\beta^{\delta\theta}\varphi^\beta_{cd}]\\
&
=-M_{\eta\theta}M_{\gamma\delta}C_\alpha^{\gamma\eta}\varphi_{ab}^\alpha
\varphi^\delta_{cf}\varphi^\theta_{df}
+M_{\gamma\delta}M_{\eta\theta}C_\alpha^{\gamma\eta}
C_\beta^{\delta\theta}\varphi^\alpha_{ab}\varphi^\beta_{cd}\\
&
=-M_{\eta\theta}M_{\gamma\delta}C_\alpha^{\gamma\eta}\varphi_{ab}^\alpha
\varphi^\delta_{cf}\varphi^\theta_{df}
+(C_\alpha^{\gamma\eta}C_\beta^{\delta\theta}
M_{\gamma\delta}M_{\eta\theta})\varphi^\alpha_{ab}\varphi^\beta_{cd}
\end{align*}
which implies
$$
M_{\gamma\delta}\varphi_{cf}^\delta
M_{\eta\theta}\varphi^\theta_{df}C_\alpha^{\gamma\eta}\varphi_{ab}^\alpha
=\frac{1}{2}(C_\alpha^{\gamma\eta}C_\beta^{\delta\theta}
M_{\gamma\delta}M_{\eta\theta})\varphi^\alpha_{ab}\varphi^\beta_{cd}.
$$
Then we have \be 2(B_{acbd}-B_{adbc})=
(C_\alpha^{\gamma\eta}C_\beta^{\delta\theta}
M_{\gamma\delta}M_{\eta\theta})\varphi^\alpha_{ab}
\varphi^\beta_{cd}. %%\eqno(1.3.8)
\ee

Therefore, combining (1.3.7) and (1.3.8), we can reformulate the
curvature evolution equation (1.3.4) as follows.

%%\vskip 0.2cm\noindent{\bf Proposition 1.3.4} (Hamilton
%%\cite{Ha86}) \emph{
% CZ-BASELINE source-line:1660 source-kind:proposition source-id:1.3.4
\begin{proposition}[evolution of the curvature operator]
  \label{prop:cz-evolution-of-the-curvature-operator}
  \dcref{Proposition 1.3.4}
  \group{Cao--Zhu Evolution of Curvatures}
  \source{cao-zhu}
Let
$R_{abcd}=M_{\alpha\beta}\varphi_{ab}^\alpha\varphi^\beta_{cd}.$
Then under the Ricci flow $(1.1.5),$ $M_{\alpha\beta}$ satisfies
the evolution equation \be \frac{\partial
M_{\alpha\beta}}{\partial t} =\Delta
M_{\alpha\beta}+M_{\alpha\beta}^2+M_{\alpha\beta}^\# %\eqno(1.3.9)
\ee where $M_{\alpha\beta}^2=M_{\alpha\gamma}M_{\beta\gamma}$ is
the operator square and
$M_{\alpha\beta}^\#=(C_\alpha^{\gamma\eta}C_\beta^{\delta\theta}
M_{\gamma\delta}M_{\eta\theta})$ is the Lie algebra square.
%%$$\eqno \#$$ \vskip 0.3cm
\end{proposition}

Let us now consider the operator $M_{\alpha\beta}^\#$ in
dimensions 3 and 4 in more detail.

In dimension $3$, let $\omega_1,\omega_2,\omega_3$ be a positively
oriented orthonormal basis for one-forms. Then
$$
\varphi^1=\sqrt{2}\omega_1\wedge\omega_2,\qquad\varphi^2
=\sqrt{2}\omega_2\wedge\omega_3,
\qquad\varphi^3=\sqrt{2}\omega_3\wedge\omega_1
$$
form an orthonormal basis for two-forms $\Lambda^2$. Write
$\varphi ^{\alpha} = \{\varphi ^{\alpha}_{ab}\}, \alpha =1,2,3,$
as
\begin{gather*}
(\varphi^1_{ab})= \begin{pmatrix}
0& \frac{\sqrt{2}}{2}& 0\\
-\frac{\sqrt{2}}{2}& 0&  0\\
0&  0&  0\end{pmatrix},\qquad (\varphi^2_{ab})=\begin{pmatrix} 0&
0& 0\\0& 0& \frac{\sqrt{2}}{2}\\0& -\frac{\sqrt{2}}{2}&
0\end{pmatrix},\\
(\varphi^3_{ab})=\begin{pmatrix} 0& 0& -\frac{\sqrt{2}}{2}\\0& 0&
0\\\frac{\sqrt{2}}{2}& 0& 0\end{pmatrix},
\end{gather*}
then
$$\arraycolsep=1.5pt\begin{array}{rcl}
[\varphi^1,\varphi^2]&  =&  \left(\begin{array}{ccc} 0&
\frac{\sqrt{2}}{2}& 0\\
-\frac{\sqrt{2}}{2}& 0& 0\\0& 0&
0\end{array}\right)\left(\begin{array}{ccc} 0& 0& 0\\0& 0&
-\frac{\sqrt{2}}{2}\\0& \frac{\sqrt{2}}{2}&
0\end{array}\right)-\left(\begin{array}{ccc} 0& 0& 0\\0& 0&
\frac{\sqrt{2}}{2}\\0& -\frac{\sqrt{2}}{2}&
0\end{array}\right)\left(\begin{array}{ccc}
0&  -\frac{\sqrt{2}}{2}&  0\\
\frac{\sqrt{2}}{2}&  0&  0\\
0&  0&  0\end{array}\right)\\[7mm]
&  =&  \left(\begin{array}{ccc}
0&  0&  -\frac{1}{2}\\0&  0&  0\\
\frac{1}{2}&  0&  0\end{array}\right)\\[7mm]
&  =& {\displaystyle \frac{\sqrt{2}}{2}\varphi^3.}
\end{array}
$$
So
$C^{123}=\<[\varphi^1,\varphi^2],\varphi^3\>=\frac{\sqrt{2}}{2}$,
in particular
$$
C^{\alpha\beta\gamma}=\begin{cases} \pm\frac{\sqrt{2}}{2},&
\text{ if }\;
\alpha\neq\beta\neq\gamma,\\
0,&  \text{ otherwise}.
\end{cases}
$$
Hence the matrix $M^\#=(M^{\#}_{\alpha\beta})$ is just the adjoint
matrix of $M=(M_{\alpha\beta})$: \be
M^\#=\det M\cdot {}^tM^{-1}. %\eqno(1.3.10)
\ee

In dimension 4, we can use the Hodge star operator $\*$ to
decompose the space of two-forms $\Lambda^2$ as
$$
\Lambda^2={\Lambda^2_+}{\oplus}{\Lambda^2_-}
$$
where $\Lambda^2_+$ (resp. $\Lambda^2_-$) is the eigenspace of the
star operator with eigenvalue $+1$ (resp. $-1$). Let $\omega_1,
\omega_2, \omega_3, \omega_4$ be a positively oriented orthonormal
basis for one-forms. A basis for $\Lambda^2_+$ is then given by
$$
\varphi^1=\omega_1\wedge\omega_2 +
\omega_3\wedge\omega_4,\quad\varphi^2 =\omega_1\wedge\omega_3 +
\omega_4\wedge\omega_2, \quad\varphi^3=\omega_1\wedge\omega_4 +
\omega_2\wedge\omega_3,
$$
while a basis for $\Lambda^2_-$ is given by
$$
\psi^1=\omega_1\wedge\omega_2 - \omega_3\wedge\omega_4,\quad\psi^2
=\omega_1\wedge\omega_3 - \omega_4\wedge\omega_2,
\quad\psi^3=\omega_1\wedge\omega_4 - \omega_2\wedge\omega_3.
$$
In particular, $\{\varphi^1, \varphi^2, \varphi^3, \psi^1, \psi^2,
\psi^3 \}$ forms an orthonormal basis for the space of two-forms
$\Lambda^2$. By using this basis we obtain a block decomposition
of the curvature operator matrix $M$ as
$$
M=(M_{\alpha\beta})=\left(\begin{array}{cc} A&  B\\ {}^tB&
C\end{array}\right).
$$
Here $A, B$ and $C$ are $3\times 3$ matrices with $A$ and $C$
being symmetric. Then we can write each element of the basis as a
skew-symmetric $4\times 4$ matrix and compute as above to get \be
M^\#=(M_{\alpha\beta}^\#)=2\left(\begin{array}{cc}
A^\#&  B^\#\\ {}^tB^\#&  C^\#\end{array}\right), %\eqno(1.3.11)
\ee where $A^\#,B^\#,C^\#$ are the adjoint of $3\times3$
submatrices as before.

For later applications in Chapter 5, we now give some computations
for the entries of the matrices $A$, $C$ and $B$ as follows. First
for the matrices $A$ and $C$, we have
$$
A_{11} = Rm(\varphi^1,\varphi^1) = R_{1212} + R_{3434} + 2R_{1234}
$$
$$
A_{22} = Rm(\varphi^2,\varphi^2) = R_{1313} + R_{4242} + 2R_{1342}
$$
$$
A_{33} = Rm(\varphi^3,\varphi^3) = R_{1414} + R_{2323} + 2R_{1423}
$$
and
$$
C_{11} = Rm(\psi^1,\psi^1) = R_{1212} + R_{3434} -2R_{1234}
$$
$$
C_{22} = Rm(\psi^2,\psi^2) = R_{1313} + R_{4242} -2R_{1342}
$$
$$
C_{33} = Rm(\psi^3,\psi^3) = R_{1414} + R_{2323} -2R_{1423}.
$$
By the Bianchi identity
$$
R_{1234} + R_{1342} + R_{1423} =0,
$$
so we have
$$
trA = trC = \frac{1}{2}R.
$$
Next for the entries of the matrix $B$, we have
$$
B_{11} = Rm(\varphi^1,\psi^1) = R_{1212} - R_{3434}
$$
$$
B_{22} = Rm(\varphi^2,\psi^2) = R_{1313} - R_{4242}
$$
$$
B_{33} = Rm(\varphi^3,\psi^3) = R_{1414} - R_{2323}
$$
and
$$
B_{12} = Rm(\varphi^1,\psi^2) = R_{1213} + R_{3413} -R_{1242} -
R_{3442} \ \ \mbox{ etc. }
$$
Thus the entries of $B$ can be written as
$$
B_{11} = \frac{1}{2}(R_{11} + R_{22} -R_{33} -R_{44})
$$
$$
B_{22} = \frac{1}{2}(R_{11} + R_{33} -R_{44} -R_{22})
$$
$$
B_{33} = \frac{1}{2}(R_{11} + R_{44} -R_{22} -R_{33})
$$
and
$$
B_{12} = R_{23} - R_{14} \ \ \mbox{ etc. }
$$
If we choose the frame $\{ \omega_1, \omega_2, \omega_3, \omega_4
\}$ so that the Ricci tensor is diagonal, then the matrix $B$ is
also diagonal. In particular, the matrix $B$ is identically zero
when the four-manifold is Einstein.


\section{Derivative Estimates}

In the previous section we have seen that the curvatures satisfy
nonlinear heat equations with quadratic growth terms. The
parabolic nature will give us a bound on the derivatives of the
curvatures at any time $t>0$ in terms of a bound of the
curvatures. In this section, we derive Shi's local derivative
estimate \cite{Sh89}. Our presentation follows Hamilton
\cite{Ha95F}.

We begin with the global version of the derivative estimate of Shi
\cite{Sh89}.

%%\vskip 0.2cm \noindent {\bf Theorem 1.4.1} (Shi \cite{Sh89})\emph{
% CZ-BASELINE source-line:1847 source-kind:theorem source-id:1.4.1
\begin{theorem}[global curvature derivative estimates]
  \label{thm:cz-global-curvature-derivative-estimates}
  \dcref{Theorem 1.4.1}
  \group{Cao--Zhu Derivative Estimates}
  \source{cao-zhu}
There exist constants $C_m$, $m=1,2,\ldots,$ such that if the
curvature of a complete solution to Ricci flow is bounded by
$$
|R_{ijkl}|\leq M
$$
up to time t with $0<t\leq \frac{1}{M}$, then the covariant
derivative of the curvature is bounded by
$$
|\nabla R_{ijkl}|\leq C_1M/\sqrt{t}
$$
and the $m^{th}$ covariant derivative of the curvature is bounded
by
$$
|\nabla^mR_{ijkl}|\leq C_mM/t^{\frac{m}{2}}.
$$
Here the norms are taken with respect to the evolving metric.
\end{theorem}

%%\vskip 0.1cm \noindent {\bf Proof.}
\begin{proof}
  \uses{prop:cz-evolution-of-the-riemann-curvature-tensor}
We shall only give the proof for the compact case.  The noncompact
case can be deduced from the next local derivative estimate
theorem. Let us denote the curvature tensor by $Rm$ and denote by
$A*B$ any tensor product of two tensors $A$ and $B$ when we do not
need the precise expression. We have from Proposition 1.3.1 that
\be
\frac{\partial}{\partial t}Rm=\Delta Rm+Rm*Rm.  %\eqno (1.4.1)
\ee Since
\begin{align*}
\frac{\partial}{\partial t}\Gamma^i_{jk} &  =
\frac{1}{2}g^{il}\left\{\nabla_j\frac{\partial g_{kl}}{\partial t}
+\nabla_k\frac{\partial g_{jl}}{\partial t}
-\nabla_l\frac{\partial g_{jk}}{\partial t}\right\}\\
&  = \nabla Rm ,
\end{align*}
it follows that \be \frac{\partial}{\partial t}(\nabla
Rm)=\Delta(\nabla Rm)
+Rm*(\nabla Rm).  %\eqno (1.4.2)
\ee Thus
\begin{align*}
\frac{\partial}{\partial t}|Rm|^2
&\leq\Delta|Rm|^2-2|\nabla Rm|^2+C|Rm|^3,\\
\frac{\partial}{\partial t}|\nabla Rm|^2 &\leq\Delta|\nabla
Rm|^2-2|\nabla^2 Rm|^2+C|Rm|\cdot|\nabla Rm|^2,
\end{align*}
for some constant $C$ depending only on the dimension $n$.

Let $A>$0 be a constant (to be determined) and set
$$
F=t|\nabla Rm|^2+A|Rm|^2.
$$
We compute
\begin{align*}
\frac{\partial F}{\partial t} &  = |\nabla
Rm|^2+t\frac{\partial}{\partial t}|\nabla Rm|^2
+A\frac{\partial}{\partial t}|Rm|^2\\
& \leq \Delta(t|\nabla Rm|^2+A|Rm|^2)+|\nabla Rm|^2(1+tC|Rm|-2A)
+CA|Rm|^3.
\end{align*}
Taking $A\geq C+1,$ we get
$$
\frac{\partial F}{\partial t}\leq \Delta F+\bar{C}M^3
$$
for some constant $\bar{C}$ depending only on the dimension $n$.
We then obtain
$$
F\leq F(0)+\bar{C}M^3t\leq (A+\bar{C})M^2,
$$
and then
$$
|\nabla Rm|^2\leq(A+\bar{C})M^2/t.
$$

The general case follows in the same way. If we have bounds
$$
|\nabla^k Rm|\leq C_kM/t^\frac{k}{2},
$$
we know from (1.4.1) and (1.4.2) that
$$
\frac{\partial}{\partial t}|\nabla^kRm|^2
\leq\Delta|\nabla^kRm|^2-2|\nabla^{k+1}Rm|^2+\frac{CM^3}{t^k},
$$
and
$$
\frac{\partial}{\partial t}|\nabla^{k+1}Rm|^2
\leq\Delta|\nabla^{k+1}Rm|^2-2|\nabla^{k+2}Rm|^2
+CM|\nabla^{k+1}Rm|^2+\frac{CM^3}{t^{k+1}}.
$$
Let $A_k > 0$ be a constant (to be determined) and set
$$
F_k=t^{k+2}|\nabla^{k+1}Rm|^2+A_k t^{k+1}|\nabla^kRm|^2.
$$
Then
\begin{align*}
\frac{\partial }{\partial t}F_k &  =(k+2)t^{k+1}|\nabla ^{k+1}
Rm|^2
+t^{k+2}\frac{\partial}{\partial t}|\nabla ^{k+1} Rm|^2\\
&\quad  +A_k (k+1) t^k |\nabla ^{k} Rm|^2
+ A_k t^{k+1}\frac{\partial}{\partial t}|\nabla ^{k}Rm|^2\\
& \leq (k+2)t^{k+1}|\nabla ^{k+1} Rm|^2 \\
&\quad+t^{k+2}\left[\Delta|\nabla^{k+1}Rm|^2 -2|\nabla^{k+2}Rm|^2
+CM|\nabla^{k+1}Rm|^2+\frac{CM^3}{t^{k+1}} \right]\\
&\quad  +A_k (k+1) t^k |\nabla ^{k} Rm|^2  \\
&\quad+ A_k t^{k+1}\left[\Delta|\nabla^kRm|^2
-2|\nabla^{k+1}Rm|^2+\frac{CM^3}{t^k} \right]\\
& \leq \Delta F_k + C_{k+1}M^2
\end{align*}
for some positive constant $C_{k+1}$, by choosing $A_k$ large
enough. This implies that
$$
|\nabla^{k+1}Rm|\leq\frac{C_{k+1}M}{t^\frac{k+1}{2}}.
$$
%%$$\eqno\#$$ \vskip 0.3cm
\end{proof}

The above derivative estimate is a somewhat standard Bernstein
estimate in PDEs. By using a cutoff argument, we will derive the
following local version in \cite{Sh89}, which is called
\textbf{Shi's derivative estimate}\index{Shi's derivative
estimate}. The proof is adapted from Hamilton \cite{Ha95F}.

%%\vskip 0.2cm \noindent {\bf Theorem 1.4.2 (Shi \cite{Sh89})}\emph{
% CZ-BASELINE source-line:1970 source-kind:theorem source-id:1.4.2
\begin{theorem}[local curvature derivative estimates]
  \label{thm:cz-local-curvature-derivative-estimates}
  \dcref{Theorem 1.4.2}
  \group{Cao--Zhu Derivative Estimates}
  \source{cao-zhu}
There exist positive constants $\theta, C_k, k=1, 2, \ldots,$
depending only on the dimension with the following property.
Suppose that the curvature of a solution to the Ricci flow is
bounded
$$
|Rm|\leq M, \ \ \ \ \mbox{on} \
U\times\left[0,\frac{\theta}{M}\right]
$$
where $U$ is an open set of the manifold. Assume that the closed
ball ${B_0(p,r)}$, centered at $p$ of radius $r$ with respect to
the metric at $t=0$, is contained in $U$ and the time $t\leq
{\theta}/{M}$. Then we can estimate the covariant derivatives of
the curvature at $(p,t)$ by
$$
|\nabla Rm(p,t)|^2\leq C_1M^2\(\frac{1}{r^2}+\frac{1}{t}+M\),
$$
and the $k^{th}$ covariant derivative of the curvature at $(p,t)$
by
$$
|\nabla ^k Rm(p,t)|^2\leq
C_kM^2\(\frac{1}{r^{2k}}+\frac{1}{t^k}+M^k\).
$$
\end{theorem}

%%\vskip 0.1cm \noindent {\bf Proof.}
\begin{proof}
Without loss of generality, we may assume $r\leq\theta/\sqrt{M}$
and the exponential map at $p$ at time $t=0$ is injective on the
ball of radius $r$ (by passing to a local cover if necessary, and
pulling back the local solution of the Ricci flow to the ball of
radius $r$ in the tangent space at $p$ at time $t=0$).

Recall
\begin{align*}
\frac{\partial}{\partial t}|Rm|^2
&\leq\Delta|Rm|^2-2|\nabla Rm|^2+C|Rm|^3, \\
\frac{\partial}{\partial t}|\nabla Rm|^2 &\leq\Delta|\nabla
Rm|^2-2|\nabla^2 Rm|^2+C|Rm|\cdot|\nabla Rm|^2.
\end{align*}
Define
$$
S=(BM^2+|Rm|^2)|\nabla Rm|^2
$$
where $B$ is a positive constant to be determined. By choosing
$B\geq {C^2}/{4}$ and using the Cauchy inequality, we have
\begin{align*}
\frac{\partial}{\partial t}S
&  \leq \Delta S-2BM^2|\nabla^2Rm|^2-2|\nabla Rm|^4 \\
&\quad+CM|\nabla Rm|^2\cdot|\nabla^2Rm|+CBM^3|\nabla Rm|^2\\
&  \leq \Delta S-|\nabla Rm|^4+CB^2M^6\\
&  \leq \Delta S-\frac{S^2}{(B+1)^2M^4}+CB^2M^6.
\end{align*}
If we take
$$
F=b(BM^2+|Rm|^2)|\nabla Rm|^2/M^4=bS/M^4,
$$
and $b\leq \min\{1/(B+1)^2,1/CB^2\}$, we get \be
\frac{\partial F}{\partial t}\leq\Delta F-F^2+M^2.%\eqno(1.4.3)
\ee

We now want to choose a cutoff function $\varphi$ with the support
in the ball $B_0(p,r)$ such that at $t=0$,
$$
\varphi(p)=r,\ \ \ 0\leq\varphi\leq Ar,
$$
and
$$ |\nabla\varphi|\leq A,\ \ \
|\nabla^2\varphi|\leq\frac{A}{r}
$$
for some positive constant $A$ depending only on the dimension.
Indeed, let $g:(-\infty,+\infty)\rightarrow[0,+\infty)$ be a
smooth, nonnegative function satisfying
$$
g(u)= \begin{cases}
         1, & u\in(-\frac{1}{2},\frac{1}{2}),\\
         0, & \text{outside }\; (-1,1).
       \end{cases}
$$
Set
$$
\varphi=rg\(\frac{s^2}{r^2}\),
$$
where $s$ is the geodesic distance function from $p$ with respect
to the metric at $t=0$. Then
$$
\nabla\varphi=\frac{1}{r}g'\(\frac{s^2}{r^2}\)\cdot 2s\nabla s
$$
and hence $$|\nabla\varphi|\leq 2C_1.$$ Also,
$$
\nabla^2\varphi=\frac{1}{r}g''\(\frac{s^2}{r^2}\)\frac{1}{r^2}4s^2\nabla
s\cdot\nabla s+\frac{1}{r}g'\(\frac{s^2}{r^2}\)2\nabla s\cdot
\nabla s+\frac{1}{r}g'\(\frac{s^2}{r^2}\)\cdot 2s\nabla^2 s.
$$
Thus, by using the standard Hessian comparison,
\begin{align*}
|\nabla^2\varphi|&  \leq \frac{C_1}{r}+\frac{C_1}{r}|s\nabla^2s|\\
&  \leq \frac{C_1}{r}\(1+s\(\frac{C_2}{s}+\sqrt{M}\)\)\\
&  \leq \frac{C_3}{r}.
\end{align*}
Here $C_1, C_2$ and $C_3$ are positive constants depending only on
the dimension.

Now extend $\varphi$ to $U\times[0,\frac{\theta}{M}]$ by letting
$\varphi$ to be zero outside $B_0(p,r)$ and independent of time.
Introduce the barrier function \be
H=\frac{(12+4\sqrt{n})A^2}{\varphi^2}+\frac{1}{t}+M %%NEW
\ee which is defined and smooth on the set
$\{\varphi>0\}\times(0,T]$.

As the metric evolves, we will still have $0\leq\varphi\leq Ar$
(since $\varphi$ is independent of time $t$); but
$|\nabla\varphi|^2$ and $\varphi|\nabla^2\varphi|$ may increase.
By continuity it will be a while before they double.

\end{proof}

%%\vskip 0.2cm\noindent{\bf Claim 1}: \emph{
% CZ-BASELINE source-line:2086 source-kind:claim source-id:claim-1
\begin{proposition}[cutoff barrier is a strict supersolution]
  \label{prop:cz-cutoff-barrier-is-a-strict-supersolution}
  \dcref{Claim 1}
  \group{Cao--Zhu Derivative Estimates}
  \source{cao-zhu}
As long as
$$
|\nabla\varphi|^2\leq 2A^2,\ \ \varphi|\nabla^2\varphi|\leq 2A^2,
$$
we have
$$
\frac{\partial H}{\partial t}>\Delta H-H^2+M^2.$$
\end{proposition}

\begin{proof}
Indeed, by the definition of $H$, we have
\begin{align*}
H^2 &>\frac{(12+4\sqrt{n})^2A^4}{\varphi^4}+\frac{1}{t^2}+M^2,\\
\frac{\partial H}{\partial t}&=-\frac{1}{t^2},
\end{align*}
and
\begin{align*}
\Delta H&  = (12+4\sqrt{n})A^2\Delta \(\frac{1}{\varphi^2}\)\\
&  = (12+4\sqrt{n})A^2\(\frac{6|\nabla\varphi|^2
-2\varphi\Delta\varphi}{\varphi^4}\)\\
&  \leq (12+4\sqrt{n})A^2\(\frac{12A^2+4\sqrt{n}A^2}{\varphi^4}\)\\
&  = \frac{(12+4\sqrt{n})^2A^4}{\varphi^4}.
\end{align*}
Therefore,
$$
H^2>\Delta H-\frac{\partial H}{\partial t}+M^2.
$$
\end{proof}

%%\vskip 0.2cm\noindent {\bf Claim 2}:\emph{\
% CZ-BASELINE source-line:2115 source-kind:claim source-id:claim-2
\begin{proposition}[short-time control of cutoff derivatives]
  \label{prop:cz-short-time-control-of-cutoff-derivatives}
  \dcref{Claim 2}
  \group{Cao--Zhu Derivative Estimates}
  \source{cao-zhu}
If the constant $\theta>0$ is small enough compared to $b$, $B$
and $A$, then we have the following property: as long as $r\leq
{\theta}/{\sqrt{M}}$, $t\leq {\theta}/{M}$ and $F\leq H$, we will
have
$$
|\nabla \varphi|^2\le 2A^2\quad \text{ and }\quad
\varphi|\nabla^2\varphi|\le 2A^2.
$$
\end{proposition}

\begin{proof}
Indeed, by considering the evolution of $\nabla\varphi$, we have
\begin{align*}
\frac{\partial}{\partial t}\nabla_a\varphi
& =\frac{\partial}{\partial t}(F^i_a\nabla_i\varphi)\\
&  = F^i_a\nabla_i\(\frac{\partial \varphi}{\partial t}\)
+\nabla_i\varphi R_k^iF_a^k\\
&  = R_{ab}\nabla_b\varphi
\end{align*}
which implies
$$
\frac{\partial}{\partial t}|\nabla \varphi|^2\le CM|\nabla
\varphi|^2,
$$
and then
$$|\nabla \varphi|^2\le A^2e^{CMt}\le 2A^2,
$$
provided $t\le \theta/M$ with $\theta\le \log 2/C$.

By considering the evolution of $\nabla^2\varphi$, we have
\begin{align*}
\frac{\partial}{\partial t}(\nabla_a\nabla_b\varphi)
& =\frac{\partial}{\partial t}(F_a^iF_b^j\nabla_i\nabla_j\varphi)\\
&  = \frac{\partial}{\partial t}\(F_a^iF_b^j
\(\frac{\partial^2\varphi}{\partial x^i\partial x^j}
-\Gamma_{ij}^k\frac{\partial\varphi}{\partial x^k}\)\)\\
&  = \nabla_a\nabla_b\(\frac{\partial\varphi}{\partial t}\)
+R_{ac}\nabla_b\nabla_c\varphi+R_{bc}\nabla_a\nabla_c\varphi\\
&\quad
+(\nabla_cR_{ab}-\nabla_aR_{bc}-\nabla_bR_{ac})\nabla_c\varphi
\end{align*}
which implies
%%\setcounter{equation}{4}%%%%1.4.4 is missing
\be \frac{\partial}{\partial t}|\nabla^2\varphi| \le
C|Rm|\cdot|\nabla^2\varphi|
+C|\nabla Rm|\cdot|\nabla \varphi|.%\eqno(1.4.5)
\ee By assumption $F\le H$, we have \be |\nabla Rm|^2\le
\frac{2M^2}{bB}\(\frac{(12+4\sqrt{n})A^2}{\varphi^2}+\frac{1}{t}\),\quad
\mbox{for}\quad t\le \theta/M. %\eqno(1.4.6)
\ee Thus by noting $\varphi$ independent of $t$ and $\varphi\le
Ar$, we get from (1.4.5) and (1.4.6) that
$$
\frac{\partial}{\partial t}(\varphi|\nabla^2\varphi|) \le
CM\(\varphi|\nabla^2\varphi|+1+\frac{r}{\sqrt{t}}\)
$$
which implies
\begin{align*}
\varphi|\nabla^2\varphi| &  \le
e^{CMt}\left[(\varphi|\nabla^2\varphi|)|_{t=0}
+CM\int_0^t\(1+\frac{r}{\sqrt{t}}\)dt\right]\\
&  \le e^{CMt}\left[A^2+CM(t+2r\sqrt{t})\right]\\
&  \le 2A^2
\end{align*}
provided $r\le\theta/\sqrt{M}$, and $t\le\theta/M$ with $\theta$
small enough. Therefore we have obtained Claim 2.
\end{proof}

\begin{proof}
\proves{thm:cz-local-curvature-derivative-estimates}
The combination of Claim 1 and Claim 2 gives us
$$
\frac{\partial H}{\partial t}>\Delta H-H^2+M^2
$$
as long as $r\le\theta/\sqrt{M},\quad t\le\theta/M$ and $F\le H$.
And (1.4.3) tells us
$$
\frac{\partial F}{\partial t}\le\Delta F-F^2+M^2.
$$
Then the standard maximum principle immediately gives the estimate
$$
|\nabla Rm|^2\le CM^2\(\frac{1}{\varphi^2}+\frac{1}{t}+M\)\qquad
\mbox{on }\;\{\varphi>0\} \times\Big(0,\frac{\theta}{M}\Big],
$$
which implies the first order derivative estimate.

The higher order derivative estimates can be obtained in the same
way by induction. Suppose we have the bounds
$$
|\nabla ^k Rm|^2\leq C_kM^2\(\frac{1}{r^{2k}}+\frac{1}{t^k}+M^k\).
$$
As before, by (1.4.1) and (1.4.2), we have
$$
\frac{\partial}{\partial t}|\nabla^kRm|^2
\leq\Delta|\nabla^kRm|^2-2|\nabla^{k+1}Rm|^2
+CM^3\(\frac{1}{r^{2k}}+\frac{1}{t^k}+M^k\),
$$
and
\begin{align*}
\frac{\partial}{\partial t}|\nabla^{k+1}Rm|^2
&\leq\Delta|\nabla^{k+1}Rm|^2-2|\nabla^{k+2}Rm|^2 \\
&\quad+CM|\nabla^{k+1}Rm|^2
+CM^3\(\frac{1}{r^{2(k+1)}}+\frac{1}{t^{k+1}}+M^{k+1}\).
\end{align*}
Here and in the following we denote by $C$ various positive
constants depending only on $C_k$ and the dimension.

Define
$$
S_k = \left[B_kM^2\(\frac{1}{r^{2k}}+\frac{1}{t^k}+M^k\) +
|\nabla^kRm|^2\right] \cdot |\nabla^{k+1}Rm|^2
$$
where $B_k$ is a positive constant to be determined. By choosing
$B_k$ large enough and Cauchy inequality, we have
{\allowdisplaybreaks
\begin{align*}
\frac{\partial}{\partial t}S_k &  \leq \bigg[-\frac{k}{t^{k+1}}
+\Delta|\nabla^kRm|^2-2|\nabla^{k+1}Rm|^2 \\
&\qquad+CM^3\(\frac{1}{r^{2k}}+\frac{1}{t^k}+M^k\)\bigg]
\cdot|\nabla^{k+1}Rm|^2\\
&\quad+\left[B_kM^2\(\frac{1}{r^{2k}}+\frac{1}{t^k}+M^k\)
+|\nabla^{k}Rm|^2\right] \\
&\qquad\cdot\bigg[\Delta|\nabla^{k+1}Rm|^2
-2|\nabla^{k+2}Rm|^2 +CM|\nabla^{k+1}Rm|^2 \\
&\qquad+CM^3\(\frac{1}{r^{2(k+1)}}+\frac{1}{t^{k+1}}+M^{k+1}\)\bigg]\\
&  \leq \Delta S_k +
8|\nabla^{k}Rm|\cdot|\nabla^{k+1}Rm|^2\cdot|\nabla^{k+2}Rm| -
\frac{k}{t^{k+1}}|\nabla^{k+1}Rm|^2 \\
&\quad- 2|\nabla^{k+1}Rm|^4
+ CM^3|\nabla^{k+1}Rm|^2\(\frac{1}{r^{2k}}+\frac{1}{t^k}+M^k\) \\
&\quad-2|\nabla^{k+2}Rm|^2
\left[B_kM^2\(\frac{1}{r^{2k}}+\frac{1}{t^k}+M^k\)
+|\nabla^kRm|^2\right]\\
&\quad +CM|\nabla^{k+1}Rm|^2\left[B_kM^2\(\frac{1}{r^{2k}}
+\frac{1}{t^k}+M^k\)+|\nabla^kRm|^2\right]\\
&\quad+CM^3\(\frac{1}{r^{2(k+1)}}+\frac{1}{t^{k+1}}
+M^{k+1}\)\\
&\qquad\cdot\left[B_kM^2\(\frac{1}{r^{2k}}+\frac{1}{t^k}+M^k\)
+ |\nabla^kRm|^2\right]\\
& \leq \Delta S_k - |\nabla^{k+1}Rm|^4 +
CB_k^2M^6\(\frac{1}{r^{4k}}+\frac{1}{t^{2k}}+M^{2k}\) \\
&\quad +CB_kM^5\(\frac{1}{r^{2(2k+1)}}+\frac{1}{t^{2k+1}}+M^{2k+1}\)\\
&  \leq \Delta S_k - |\nabla^{k+1}Rm|^4 +
CB_k^2M^5\(\frac{1}{r^{2(2k+1)}}+\frac{1}{t^{2k+1}}+M^{2k+1}\)\\
&  \leq \Delta S_k -
\frac{S_k}{(B+1)^2M^4\(\frac{1}{r^{2k}}+\frac{1}{t^k}+M^k\)^2} \\
&\quad
+CB_k^2M^5\(\frac{1}{r^{2(2k+1)}}+\frac{1}{t^{2k+1}}+M^{2k+1}\).
\end{align*}
} Let $u = {1}/{r^2} + {1}/{t} +M $ and set $F_k = bS_k/u^k$. Then
\begin{align*}
\frac{\partial F_k}{\partial t} &  \leq \Delta F_k
-\frac{F_k^2}{b(B_k+1)^2M^4u^k}
+ bCB_k^2M^5u^{k+1} + kF_ku\\
&  \leq \Delta F_k -\frac{F_k^2}{2b(B_k+1)^2M^4u^k} +
b(C+2k^2)(B_k+1)^2M^4u^{k+2}.
\end{align*}
By choosing $b \leq 1/(2(C+2k^2)(B_k+1)^2M^4)$, we get
$$
\frac{\partial F_k}{\partial t}\leq\Delta F_k - \frac{1}{u^k}F_k^2
+ u^{k+2}.
$$
Introduce
$$
H_k = 5(k+1)(2(k+1)+1+\sqrt{n})A^2\varphi^{-2(k+1)} + Lt^{-(k+1)}
+ M^{k+1},
$$
where $L \geq k+2$. Then by using Claim 1 and Claim 2, we have
$$
\frac{\partial H_k}{\partial t} = -(k+1)Lt^{-(k+2)},
$$
$$
\Delta H_k \leq 20(k+1)^2(2(k+1)+1+\sqrt{n})A^4\varphi^{-2(k+2)}
$$
and
$$
H_k^2 > 25(k+1)^2(2(k+1)+1+\sqrt{n})A^4\varphi^{-4(k+1)} +
L^2t^{-2(k+1)} + M^{2(k+1)}.
$$
These imply
$$
\frac{\partial H_k}{\partial t} >\Delta H_k - \frac{1}{u^k}H_k^2 +
u^{k+2}.
$$
Then the maximum principle immediately gives the estimate
$$
F_k \leq H_k.
$$
In particular,
\begin{align*}
&\frac{b}{u^k}B_kM^2\(\frac{1}{r^{2k}}
+\frac{1}{t^k}+M^k\)\cdot|\nabla^{k+1}Rm|^2 \\
& \leq 5(k+1)\(2(k+1)+1+\sqrt{n}\)A^2\varphi^{-2(k+1)} +
Lt^{-(k+1)} + M^{k+1}.
\end{align*}
So by the definition of $u$ and the choosing of $b$, we obtain the
desired estimate
$$
|\nabla ^{k+1} Rm|^2 \leq
C_{k+1}M^2\(\frac{1}{r^{2(k+1)}}+\frac{1}{t^{k+1}}+M^{k+1}\).
$$

Therefore we have completed the proof of the theorem.
%%$$\eqno \#$$\vskip 1cm
\end{proof}


\section{Variational Structure and Dynamic Property}

In this section, we introduce two functionals of Perelman
\cite{P1}, $\mathcal F$ and $\mathcal W$, and discuss their
relations with the Ricci flow. Our presentation here follows
sections 1.1-1.3 of Perelman \cite{P1}.

It was not known whether the Ricci flow is a gradient flow until
Perelman \cite{P1} showed that the Ricci flow is, in a certain
sense, the gradient flow of the functional $\mathcal F$. If we
consider the Ricci flow as a dynamical system on the space of
Riemannian metrics, then these two functionals are of Lyapunov
type for this dynamical system. Obviously, Ricci flat metrics are
fixed points of the dynamical system. When we consider the space
of Riemannian metrics modulo diffeomorphism and scaling, fixed
points of the Ricci flow dynamical system correspond to steady, or
shrinking, or expanding Ricci solitons. The following concept
corresponds to a periodic orbit.

%%\vskip 0.2cm\noindent{\bf Definition 1.5.1}\
% CZ-BASELINE source-line:2338 source-kind:definition source-id:1.5.1
\begin{definition}[Ricci breathers]
  \label{def:cz-ricci-breathers}
  \dcref{Definition 1.5.1}
  \group{Cao--Zhu Variational Structure and Dynamic Property}
  \source{cao-zhu}
A metric $g_{ij}(t)$ evolving by the Ricci flow is called a {\bf
breather}\index{breather} if for {\bf some} $t_1<t_2$ and
$\alpha>0$ the metrics $\alpha g_{ij}(t_1)$ and $g_{ij}(t_2)$
differ only by a diffeomorphism; the case
$\alpha=1,\;\alpha<1,\;\alpha>1$ correspond to {\bf steady,
shrinking and expanding
breathers}\index{breather!steady}\index{breather!shrinking}\index{breather!expanding},
respectively.
\end{definition}

Clearly,\; (steady,\; shrinking\; or\; expanding)\; Ricci\;
solitons\; are\; trivial breathers for which the metrics
$g_{ij}(t_1)$ and $g_{ij}(t_2)$ differ only by diffeomorphism and
scaling for every pair $t_1$ and $t_2$.

We always assume $M$ is a compact $n$-dimensional manifold in this
section. Let us first consider the functional \be
{\mathcal{F}} (g_{ij},f)=\int_M(R+|\nabla f|^2)e^{-f}dV %\eqno(1.5.1)
\ee of Perelman \cite{P1} defined on the space of Riemannian
metrics, and smooth functions on $M$. Here $R$ is the scalar
curvature of $g_{ij}$.

%\vskip 0.2cm\noindent {\bf Lemma 1.5.2 (Perelman \cite{P1})}\emph{
% CZ-BASELINE source-line:2362 source-kind:lemma source-id:1.5.2
\begin{lemma}[first variation of Perelman’s F-functional]
  \label{lem:cz-first-variation-of-perelmans-f-functional}
  \dcref{Lemma 1.5.2}
  \group{Cao--Zhu Variational Structure and Dynamic Property}
  \source{cao-zhu}
If $\delta g_{ij}=v_{ij}$ and $\delta f=h$ are variations of
$g_{ij}$ and $f$ respectively, then the first variation of
${\mathcal{F}}$ is given by
$$
\delta{\mathcal{F}}(v_{ij},h) =\!\!\int_M\left[-v_{ij}
(R_{ij}+\nabla_i\nabla_jf)+\(\frac{v}{2}-h\)(2\Delta f-|\nabla
f|^2+R)\right]e^{-f}dV
$$
where $v=g^{ij}v_{ij}$.
\end{lemma}

%%\vskip 0.1cm\noindent{\bf Proof.}\
\begin{proof}
In any normal coordinates at a fixed point, we have
\begin{align*}
\delta R^h_{ijl}&  = \frac{\partial}{\partial
x^i}(\delta\Gamma^h_{jl})-\frac{\partial}{\partial
x^j}(\delta\Gamma^h_{il})\\
&  = \frac{\partial}{\partial x^i}\left[\frac{1}{2}g^{hm}(\nabla_j
v_{lm}+\nabla_lv_{jm}-\nabla_mv_{jl})\right]\\
&\quad -\frac{\partial}{\partial
x^j}\left[\frac{1}{2}g^{hm}(\nabla_i
v_{lm}+\nabla_lv_{im}-\nabla_mv_{il})\right],\\
\delta R_{jl}&  = \frac{\partial}{\partial
x^i}\left[\frac{1}{2}g^{im}(\nabla_j
v_{lm}+\nabla_lv_{jm}-\nabla_mv_{jl})\right]\\
&\quad -\frac{\partial}{\partial
x^j}\left[\frac{1}{2}g^{im}(\nabla_i
v_{lm}+\nabla_lv_{im}-\nabla_mv_{il})\right]\\
&  = \frac{1}{2}\frac{\partial}{\partial
x^i}[\nabla_jv^i_l+\nabla_lv^i_j-\nabla^iv_{jl}]
-\frac{1}{2}\frac{\partial}{\partial x^j}[\nabla_lv],\\
\delta R&  = \delta(g^{jl}R_{jl})\\
&  = -v_{jl}R_{jl}+g^{jl}\delta R_{jl}\\
&  = -v_{jl}R_{jl}+\frac{1}{2}\frac{\partial}{\partial x^i}
[\nabla^lv^i_l+\nabla_lv^{il}-\nabla^iv]
-\frac{1}{2}\frac{\partial}{\partial x^j}[\nabla^jv]\\
&  = -v_{jl}R_{jl}+\nabla_i\nabla_lv_{il}-\Delta v.
\end{align*}
Thus \be \delta R(v_{ij})=-\Delta
v+\nabla_i\nabla_jv_{ij}-v_{ij}R_{ij}.  %%\eqno(1.5.2)
\ee The first variation of the functional
${\mathcal{F}}(g_{ij},f)$ is
\begin{align}
& \delta\(\int_M(R+|\nabla f|^2)e^{-f}dV\)\\
&  = \int_M\bigg([\delta
R(v_{ij})+\delta(g^{ij}\nabla_if\nabla_jf)]e^{-f}dV \nn\\
&\qquad+(R+|\nabla f|^2)
\left[-he^{-f}dV+e^{-f}\frac{v}{2}dV\right]\bigg) \nn\\
&  = \int_M\bigg[-\Delta
v+\nabla_i\nabla_jv_{ij}-R_{ij}v_{ij}-v_{ij}\nabla_if\nabla_jf \nn\\
&\qquad +2\<\nabla f,\nabla h\>
+(R+|\nabla f|^2)\(\frac{v}{2}-h\)\bigg]e^{-f}dV.\nn  %%\eqno(1.5.3)
\end{align}

On the other hand,
\begin{align*}
\int_M(\nabla_i\nabla_jv_{ij}-v_{ij}\nabla_if\nabla_jf)e^{-f}dV &=
\int_M(\nabla_if\nabla_jv_{ij}
-v_{ij}\nabla_if\nabla_jf)e^{-f}dV\\
&  = -\int_M(\nabla_i\nabla_jf)v_{ij}e^{-f}dV,\\
\int_M\!2\<\nabla f,\nabla h\>e^{-f}dV &  = -2\!\int_M\!h\Delta
fe^{-f}dV+2\!\int_M\!|\nabla f|^2he^{-f}dV,
\end{align*}
and
\begin{align*}
\int_M(-\Delta v)e^{-f}dV
& = -\int_M\<\nabla f,\nabla v\>e^{-f}dV\\
& = \int_Mv\Delta fe^{-f}dV-\int_M|\nabla f|^2ve^{-f}dV.
\end{align*}
Plugging these identities into (1.5.3) the first variation formula
follows.
%%$$\eqno\#$$ \vskip 0.3cm
\end{proof}

Now let us study the functional ${\mathcal{F}}$ when the metric
evolves under the Ricci flow and the function evolves by a
backward heat equation.

%%\vskip 0.2cm\noindent {\bf Proposition 1.5.3 (Perelman
%%\cite{P1})}\ \emph{
% CZ-BASELINE source-line:2444 source-kind:proposition source-id:1.5.3
\begin{proposition}[monotonicity of Perelman’s F-functional]
  \label{prop:cz-monotonicity-of-perelmans-f-functional}
  \dcref{Proposition 1.5.3}
  \group{Cao--Zhu Variational Structure and Dynamic Property}
  \source{cao-zhu}
Let $g_{ij}(t)$ and $f(t)$ evolve according to the coupled flow
$$\arraycolsep=1.5pt\left\{\begin{array}{rcl}
\frac{\partial g_{ij}}{\partial t}&  =&  -2R_{ij},\\[2mm]
\frac{\partial f}{\partial t}&  =&  -\Delta f+|\nabla f|^2-R.
\end{array}\right.
$$
Then
$$
\frac{d}{dt}{\mathcal{F}}(g_{ij}(t),f(t))
=2\int_M|R_{ij}+\nabla_i\nabla_j f|^2e^{-f}dV
$$
and $\int_Me^{-f}dV$ is constant. In particular
${\mathcal{F}}(g_{ij}(t),f(t))$ is nondecreasing in time and the
monotonicity is strict unless we are on a steady gradient soliton.
\end{proposition}

%%\vskip 0.1cm\noindent{\bf Proof.}\
\begin{proof}
  \uses{lem:cz-first-variation-of-perelmans-f-functional}
Under the coupled flow and using the first variation formula in
Lemma 1.5.2, we have
\begin{align*}
&\frac{d}{dt}{\mathcal{F}}(g_{ij}(t),f(t)) \\
& = \int_M\bigg[-(-2R_{ij})(R_{ij}+\nabla_i\nabla_jf)\\
&\qquad +\(\frac{1}{2}(-2R) -\frac{\partial f}{\partial
t}\)(2\Delta
f-|\nabla f|^2+R)\bigg]e^{-f}dV\\
&  = \int_M[2R_{ij}(R_{ij}+\nabla_i\nabla_jf)+(\Delta f-|\nabla
f|^2)(2\Delta f-|\nabla f|^2+R)]e^{-f}dV.
\end{align*}
Now
\begin{align*}
& \int_M(\Delta f-|\nabla f|^2)(2\Delta f-|\nabla
f|^2)e^{-f}dV\\
&  = \int_M-\nabla_if\nabla_i(2\Delta f-|\nabla f|^2)e^{-f}dV\\
& =\int_M-\nabla_if(2\nabla_j(\nabla_i\nabla_jf)-2R_{ij}\nabla_jf
-2\<\nabla f,\nabla_i\nabla f\>)e^{-f}dV\\
& =\!-2\!\int_M[(\nabla_if\nabla_jf\!-\!\nabla_i\nabla_jf)
\nabla_i\nabla_jf\!-\!R_{ij}\nabla_if\nabla_jf\!-\!\<\nabla
f,\nabla_i\nabla f\>\nabla_if]e^{-f}dV\\
&  = 2\int_M[|\nabla_i\nabla_jf|^2+R_{ij}\nabla_i
f\nabla_jf]e^{-f}dV,
\end{align*}
and
\begin{align*}
&\int_M(\Delta f-|\nabla f|^2)Re^{-f}dV \\
&  = \int_M-\nabla_if\nabla_iRe^{-f}dV\\
& =2\int_M\nabla_i\nabla_jfR_{ij}e^{-f}dV
-2\int_M\nabla_if\nabla_jfR_{ij}e^{-f}dV.
\end{align*}
Here we have used the contracted second Bianchi identity.
Therefore we obtain
\begin{align*}
& \frac{d}{dt}{\mathcal{F}}(g_{ij}(t), f(t)) \\
& = \int_M[2R_{ij}(R_{ij}+\nabla_i\nabla_jf)
+2(\nabla_i\nabla_jf)(\nabla_i\nabla_jf+R_{ij})]e^{-f}dV\\
& = 2\int_M|R_{ij}+\nabla_i\nabla_jf|^2e^{-f}dV.
\end{align*}
It remains to show $\int_Me^{-f}dV$ is a constant. Note that the
volume element $dV=\sqrt{{\det}g_{ij}}\ dx$ evolves under the
Ricci flow by
\begin{align}
\frac{\partial}{\partial t}dV&  = \frac{\partial}{\partial
t}(\sqrt{\det g_{ij}}) dx\\
&  = \frac {1}{2}\(\frac{\partial }{\partial t}
\log(\det g_{ij})\) dV \nn\\
&  =\frac{1}{2}(g^{ij}\frac{\partial}{\partial t}g_{ij})dV \nn\\
&  = -RdV.  \nn  %%\eqno (1.5.4)
\end{align}
Hence
\begin{align}
\frac{\partial}{\partial t}\(e^{-f} dV\)
& = e^{-f}\(-\frac{\partial f}{\partial t}-R\)dV \\
&  = (\Delta f-|\nabla f|^2)e^{-f}dV \nn\\
&  = -\Delta(e^{-f})dV. \nn%% \eqno(1.5.5)
\end{align}
It then follows that
$$
\frac{d}{dt}\int_Me^{-f}dV=-\int_M\Delta(e^{-f})dV=0.
$$
This finishes the proof of the proposition.
%%$$\eqno \#$$ \vskip 0.3cm
\end{proof}

Next we define the associated energy \be \lambda(g_{ij})
=\inf\left\{{\mathcal{F}}(g_{ij},f)\ |\  f\in C^{\infty}(M),
\int_Me^{-f}dV=1\right\}. %%\eqno (1.5.6)
\ee

If we set $u=e^{-f/2}$, then the functional ${\mathcal{F}}$ can be
expressed in terms of $u$ as
$$
{\mathcal{F}}=\int_M(R u^2+4|\nabla u|^2)dV,
$$
and the constraint $\int_Me^{-f}dV=1$ becomes $\int_M u^2 dV=1$.
Therefore $\lambda (g_{ij})$ is just the first eigenvalue of the
operator $-4\Delta +R$.  Let $u_0>0$ be a first eigenfunction of
the operator $-4\Delta +R$ satisfying
$$
-4\Delta u_0+Ru_0=\lambda(g_{ij})u_0.
$$
The $f_0=-2\log u_0$ is a minimizer:
$$
\lambda(g_{ij})={\mathcal{F}}(g_{ij},f_0).
$$
Note that $f_0$ satisfies the equation \be
-2\Delta f_0+|\nabla f_0|^2-R=-\lambda(g_{ij}). %\eqno (1.5.7)
\ee

Observe that the evolution equation
$$
\frac{\partial f}{\partial t}=-\Delta f+|\nabla f|^2-R
$$
can be rewritten as the following linear equation
$$
\frac{\partial}{\partial t}(e^{-f})=-\Delta(e^{-f})+R(e^{-f}).
$$
Thus we can always solve the evolution equation for $f$ backwards
in time. Suppose at $t=t_0$, the infimum $\lambda(g_{ij})$ is
achieved by some function $f_0$ with $\int_M e^{-f_0}dV=1$. We
solve the backward heat equation
$$
\arraycolsep=1.5pt\left\{\begin{array}{l}
\frac{\partial f}{\partial t}=-\Delta f+|\nabla f|^2-R\\[2mm]
f|_{t=t_0}=f_0
\end{array}\right.
$$
to obtain a solution $f(t)$ for $t\le t_0$ which satisfies
$\int_Me^{-f}dV=1$. It then follows from Proposition 1.5.3 that
$$
\lambda(g_{ij}(t))\le{\mathcal{F}}(g_{ij}(t),f(t))\le{\mathcal{F}}
(g_{ij}(t_0),f(t_0))=\lambda(g_{ij}(t_0)).
$$
Also note $\lambda(g_{ij})$ is invariant under diffeomorphism.
Thus we have proved

%%\vskip 0.2cm\noindent{\bf Corollary 1.5.4}\ \emph{
% CZ-BASELINE source-line:2581 source-kind:corollary source-id:1.5.4
\begin{corollary}[lambda monotonicity and steady breathers]
  \label{cor:cz-lambda-monotonicity-and-steady-breathers}
  \dcref{Corollary 1.5.4}
  \group{Cao--Zhu Variational Structure and Dynamic Property}
  \source{cao-zhu}
\
\begin{itemize}
\item[{\rm (i)}] $\lambda(g_{ij}(t))$ is nondecreasing along the
Ricci flow and the monotonicity is strict unless we are on a
steady gradient soliton; \item[{\rm (ii)}] A steady breather is
necessarily a steady gradient soliton.
\end{itemize}
%%$$\eqno \#$$ \vskip 0.3cm
\end{corollary}

To deal with the expanding case we consider a scale invariant
version
$$
\bar{\lambda}(g_{ij})=\lambda(g_{ij})V^{\frac{2}{n}}(g_{ij}).
$$
Here $V=Vol(g_{ij})$ denotes the volume of $M$ with respect to the
metric $g_{ij}$. \pagebreak

%%\vskip 0.2cm\noindent{\bf Corollary 1.5.5}\
% CZ-BASELINE source-line:2600 source-kind:corollary source-id:1.5.5
\begin{corollary}[scale-invariant lambda monotonicity and expanding breathers]
  \label{cor:cz-scale-invariant-lambda-monotonicity-and-expanding-breathers}
  \dcref{Corollary 1.5.5}
  \group{Cao--Zhu Variational Structure and Dynamic Property}
  \source{cao-zhu}
\
\begin{itemize}
\item[{\rm (i)}] $\bar{\lambda}(g_{ij})$ is nondecreasing along
the Ricci flow whenever it is nonpositive; moreover, the
monotonicity is strict unless we are on a gradient expanding
soliton; \item[{\rm (ii)}] An expanding breather is necessarily an
expanding gradient soliton.
\end{itemize}
\end{corollary}

%%\vskip 0.1cm\noindent{\bf Proof.}\
\begin{proof}
Let $f_0$ be a minimizer of $\lambda(g_{ij}(t))$ at $t=t_0$ and
solve the backward heat equation
$$
\frac{\partial f}{\partial t}=-\Delta f+|\nabla f|^2-R
$$
to obtain $f(t)$, $t\le t_0$, with $\int_Me^{-f(t)}dV=1$. We
compute the derivative (understood in the barrier sense) at
$t=t_0$,
\begin{align*}
& \frac{d}{dt}\bar{\lambda}(g_{ij}(t)) \\
&  \ge\frac{d}{dt}({\mathcal{F}}(g_{ij}(t),f(t)) \cdot
V^{\frac{2}{n}}(g_{ij}(t)))\\
&  = V^{\frac{2}{n}}\int_M2|R_{ij}+\nabla_i\nabla_jf|^2e^{-f}dV\\
&\quad
+\frac{2}{n}V^{\frac{2-n}{n}}\int_M(-R)dV\cdot\int_M(R+|\nabla
f|^2)e^{-f}dV\\
& =2V^{\frac{2}{n}}\bigg[\int_M\bigg|R_{ij}
+\nabla_i\nabla_jf-\frac{1}{n}(R+\Delta f)g_{ij}\bigg|^2e^{-f}dV\\
&\quad +\frac{1}{n}\int_M\!(R\!+\!\Delta f)^2e^{-f}dV\!
+\!\frac{1}{n}\(\!-\!\int_M\!(R\!+\!|\nabla f|^2)e^{-f}dV\)
\(\frac{1}{V}\int_MRdV\)\bigg],
\end{align*}
where we have used the formula (1.5.4) in the computation of
$dV/dt$.

Suppose $\lambda(g_{ij}(t_0))\le 0$, then the last term on the RHS
is given by,
\begin{align*}
& \frac{1}{n}\(-\int_M(R+|\nabla
f|^2\)e^{-f}dV)\(\frac{1}{V}\int_MRdV\)\\
&  \ge \frac{1}{n}\(-\int_M(R+|\nabla f|^2)e^{-f}dV\)
\(\int_M(R+|\nabla f|^2)e^{-f}dV\)\\
&  = -\frac{1}{n}\(\int_M(R+\Delta f)e^{-f}dV\)^2.
\end{align*}
Thus at $t=t_0$,
\begin{align}
& \frac{d}{dt}\bar{\lambda}(g_{ij}(t)) \\
&  \ge 2V^{\frac{2}{n}}
\bigg[\int_M|R_{ij}+\nabla_i\nabla_jf-\frac{1}{n}(R+\Delta
f)g_{ij}|^2e^{-f}dV \nn\\
&\quad +\frac{1}{n}\(\int_M(R+\Delta f)^2e^{-f}dV
-\(\int_M(R+\Delta f)e^{-f}dV\)^2\)\bigg]\ge 0 \nn  %\eqno(1.5.8)
\end{align}
by the Cauchy-Schwarz inequality. Thus we have proved statement
(i).

We note that on an expanding breather on $[t_1,t_2]$ with $\alpha
g_{ij}(t_1)$ and $g_{ij}(t_2)$ differ only by a diffeomorphism for
some $\alpha>1$, it would necessary have
$$
\frac{dV}{dt}>0,\; \text{ for some }\;t\in[t_1,t_2].
$$
On the other hand, for every $t$,
$$
-\frac{d}{dt}\log V=\frac{1}{V}\int_MRdV\ge \lambda(g_{ij}(t))
$$
by the definition of $\lambda(g_{ij}(t))$. It follows that on an
expanding breather on $[t_1,t_2]$,
$$
\bar{\lambda}(g_{ij}(t))
=\lambda(g_{ij}(t))V^\frac{2}{n}(g_{ij}(t))<0
$$
for some $t\in[t_1,t_2]$. Then by using statement (i), it implies
$$
\bar{\lambda}(g_{ij}(t_1))<\bar{\lambda}(g_{ij}(t_2))
$$ unless we
are on an expanding gradient soliton. We also note that
$\bar{\lambda}(g_{ij}(t))$ is invariant under diffeomorphism and
scaling which implies
$$
\bar{\lambda}(g_{ij}(t_1))=\bar{\lambda}(g_{ij}(t_2)).
$$
Therefore the breather must be an expanding gradient soliton.
%%$$\eqno \#$$ \vskip 0.3cm
\end{proof}

In particular part (ii) of Corollaries 1.5.4 and 1.5.5  imply that
all compact steady or expanding Ricci solitons are gradient ones.
Combining this fact with Proposition (1.1.1), we immediately get

%%\vskip 0.2cm\noindent{\bf Proposition 1.5.6} \ \emph {
% CZ-BASELINE source-line:2693 source-kind:proposition source-id:1.5.6
\begin{proposition}[steady and expanding breathers are Einstein]
  \label{prop:cz-steady-and-expanding-breathers-are-einstein}
  \dcref{Proposition 1.5.6}
  \group{Cao--Zhu Variational Structure and Dynamic Property}
  \source{cao-zhu}
On a compact manifold, a steady or expanding breather is necessarily
an Einstein metric.
\end{proposition}
%%$$\eqno \#$$ \vskip 0.3cm

In order to handle the shrinking case, we introduce the following
important functional, also due to Perelman \cite{P1}, \be
\mathcal{W}(g_{ij},f,\tau) =\int_M [\tau (R+|\nabla
f|^2)+f-n](4\pi
\tau)^{-\frac{n}{2}}e^{-f}dV %\eqno(1.5.9)
\ee where $g_{ij}$ is a Riemannian metric, $f$ is a smooth
function on $M$, and $\tau$ is a positive scale parameter. Clearly
the functional $\mathcal{W}$ is invariant under simultaneous
scaling of $\tau$ and $g_{ij}$ (or equivalently the parabolic
scaling), and invariant under diffeomorphism. Namely, for any
positive number $a$ and any diffeomorphism $\varphi$ \be
\mathcal{W}(a\varphi^*g_{ij},\varphi^*f,a
\tau)=\mathcal{W}(g_{ij},f,\tau).% \eqno (1.5.10)
\ee

Similar to Lemma 1.5.2, we have the following first variation
formula for $\mathcal{W}$.

%%\vskip 0.2cm\noindent {\bf Lemma 1.5.7 (Perelman \cite{P1})}\emph{
% CZ-BASELINE source-line:2718 source-kind:lemma source-id:1.5.7
\begin{lemma}[first variation of Perelman’s W-functional]
  \label{lem:cz-first-variation-of-perelmans-w-functional}
  \dcref{Lemma 1.5.7}
  \group{Cao--Zhu Variational Structure and Dynamic Property}
  \source{cao-zhu}
If $v_{ij}=\delta g_{ij},\; h=\delta f,\;\mbox{ and }\;
\eta=\delta\tau$, then
\begin{align*}
& \delta \mathcal{W}(v_{ij},h,\eta)\\
&  = \int_M-\tau v_{ij}\(R_{ij}+\nabla_i\nabla_jf
-\frac{1}{2\tau}g_{ij}\)(4\pi\tau)^{-\frac{n}{2}}e^{-f}dV\\
&\quad +\int_M\(\frac{v}{2}-h-\frac{n}{2\tau}\eta\)[\tau(R+2\Delta
f
-|\nabla f|^2)+f-n-1](4\pi\tau)^{-\frac{n}{2}}e^{-f}dV\\
&\quad +\int_M \eta\(R+|\nabla f|^2
-\frac{n}{2\tau}\)(4\pi\tau)^{-\frac{n}{2}}e^{-f}dV.
\end{align*}
Here $v=g^{ij}v_{ij}$ as before.
\end{lemma}

%%\vskip 0.1cm \noindent {\bf Proof.}
\begin{proof}
  \uses{lem:cz-first-variation-of-perelmans-f-functional}
Arguing as in the proof of Lemma 1.5.2, the first variation of the
functional $\mathcal{W}$ can be computed as follows,
\begin{align*}
&\delta W(v_{ij},h,\eta)  \\
&  = \int_M[\eta(R+|\nabla f|^2)+\tau(-\Delta
v+\nabla_i\nabla_jv_{ij}-R_{ij}v_{ij}-v_{ij}\nabla_if\nabla_jf \\
&\quad +2\<\nabla f,\nabla h\>)+h](4\pi\tau)^{-\frac{n}{2}}e^{-f}dV\\
&\quad +\int_M\left[(\tau(R+|\nabla
f|^2)+f-n)\(-\frac{n}{2}\frac{\eta}{\tau}+\frac{v}{2}-h\)\right]
(4\pi\tau)^{-\frac{n}{2}}e^{-f}dV\\
& = \int_M[\eta(R+|\nabla f|^2)+h](4\pi\tau)^{-\frac{n}{2}}e^{-f}dV\\
&\quad +\int_M[-\tau v_{ij}(R_{ij}+\nabla_i\nabla_j f)
+\tau(v-2h)(\Delta f-|\nabla f|^2)](4\pi\tau)^{-\frac{n}{2}}e^{-f}dV\\
&\quad +\int_M\left[(\tau(R+|\nabla
f|^2)+f-n)\(-\frac{n}{2}\frac{\eta}{\tau}+\frac{v}{2}-h\)\right]
(4\pi\tau)^{-\frac{n}{2}}e^{-f}dV\\
& = -\int_M\tau
v_{ij}\(R_{ij}+\nabla_i\nabla_jf-\frac{1}{2\tau}g_{ij}\)
(4\pi\tau)^{-\frac{n}{2}}e^{-f}dV\\
&\quad +\int_M\(\frac{v}{2}-h-\frac{n}{2\tau}\eta\)
[\tau(R+|\nabla f|^2) \\
&\qquad+f-n+2\tau(\Delta f-|\nabla f|^2)]
(4\pi\tau)^{-\frac{n}{2}}e^{-f}dV\\
&\quad +\int_M\left[\eta\(R+|\nabla
f|^2-\frac{n}{2\tau}\)+\(h-\frac{v}{2}+\frac{n}{2\tau}\eta\)\right]
(4\pi\tau)^{-\frac{n}{2}}e^{-f}dV\\
& = \int_M-\tau
v_{ij}\(R_{ij}+\nabla_i\nabla_jf-\frac{1}{2\tau}g_{ij}\)
(4\pi\tau)^{-\frac{n}{2}}e^{-f}dV\\
&\quad +\int_M\(\frac{v}{2}-h-\frac{n}{2\tau}\eta\)[\tau(R+2\Delta
f
-|\nabla f|^2)+f-n-1](4\pi\tau)^{-\frac{n}{2}}e^{-f}dV\\
&\quad +\int_M \eta\(R+|\nabla f|^2
-\frac{n}{2\tau}\)(4\pi\tau)^{-\frac{n}{2}}e^{-f}dV.
\end{align*}
%%$$\eqno \#$$ \vskip 0.3cm
\end{proof}

The following result is analogous to Proposition 1.5.3.

%%\vskip 0.2cm\noindent {\bf Proposition 1.5.8} (Perelman \cite{P1})
%%\ \ \emph{
% CZ-BASELINE source-line:2778 source-kind:proposition source-id:1.5.8
\begin{proposition}[monotonicity of Perelman’s W-functional]
  \label{prop:cz-monotonicity-of-perelmans-w-functional}
  \dcref{Proposition 1.5.8}
  \group{Cao--Zhu Variational Structure and Dynamic Property}
  \source{cao-zhu}
If $g_{ij}(t), f(t)$ and $\tau(t)$ evolve according to the system
$$
      \left\{
       \begin{array}{lll}
       \displaystyle
\frac{\partial g_{ij}}{\partial t}=-2R_{ij},
          \\[4mm]
      \displaystyle
\frac{\partial f}{\partial t}=-\Delta f+|\nabla
f|^2-R+\frac{n}{2\tau},
          \\[4mm]
      \displaystyle
\frac{\partial \tau}{\partial t}=-1,
       \end{array}
      \right.
$$
then we have the identity
$$
\frac{d} {d t} \mathcal{W}(g_{ij}(t),f(t),\tau(t))=\int_M
2\tau\left|R_{ij}+\nabla_i\nabla_jf-\frac
{1}{2\tau}g_{ij}\right|^2 (4\pi\tau)^{-\frac{n}{2}} e^{-f}dV
$$
and $\int_M(4\pi\tau)^{-\frac{n}{2}}e^{-f}dV$ is constant. In
particular $\mathcal{W}(g_{ij}(t),f(t),\tau(t))$ is nondecreasing
in time and the monotonicity is strict unless we are on a
shrinking gradient soliton.
\end{proposition}

%%\vskip 0.1cm \noindent {\bf Proof.}
\begin{proof}
  \uses{lem:cz-first-variation-of-perelmans-w-functional}
Using Lemma 1.5.7, we have
\begin{align}
&\frac{d} {d t} \mathcal{W}(g_{ij}(t),f(t),\tau(t)) \\
&  = \int_M 2\tau R_{ij}\(R_{ij}+\nabla_i\nabla_jf
-\frac{1}{2\tau}g_{ij}\)(4\pi\tau)^{-\frac{n}{2}}e^{-f}dV \nn\\
&\quad +\int_M(\Delta f-|\nabla f|^2)[\tau(R+2\Delta f-|\nabla
f|^2)+f](4\pi\tau)^{-\frac{n}{2}}e^{-f}dV \nn\\
&\quad -\int_M\(R+|\nabla f|^2-\frac{n}{2\tau}\)
(4\pi\tau)^{-\frac{n}{2}}e^{-f}dV. \nn  %%\eqno (1.5.11)
\end{align}
Here we have used the fact that $\int_M(\Delta f-|\nabla
f|^2)e^{-f}dV=0.$

The second term on the RHS of (1.5.11) is {\allowdisplaybreaks
\begin{align*}
& \int_M(\Delta f-|\nabla f|^2)[\tau(R+2\Delta f-|\nabla f|^2)+f]
(4\pi\tau)^{-\frac{n}{2}}e^{-f}dV\\
& = \int_M(\Delta f-|\nabla f|^2)(2\tau\Delta f-\tau|\nabla
f|^2)(4\pi\tau)^{-\frac{n}{2}}e^{-f}dV \\
&\quad-\int_M|\nabla f|^2(4\pi\tau)^{-\frac{n}{2}}e^{-f}dV
 +\tau\int_M(-\nabla_if)(\nabla_iR)
(4\pi\tau)^{-\frac{n}{2}}e^{-f}dV\\
& = \tau\int_M(-\nabla_if)(\nabla_i(2\Delta f-|\nabla
f|^2))(4\pi\tau)^{-\frac{n}{2}}e^{-f}dV \\
&\quad-\int_M\Delta f(4\pi\tau)^{-\frac{n}{2}}e^{-f}dV
-2\tau\int_M\nabla_if\nabla_jR_{ij}(4\pi\tau)^{-\frac{n}{2}}
e^{-f}dV\\
& = -2\tau \int_M(\nabla_if)(\nabla_i\Delta f-\<\nabla
f,\nabla_i\nabla f\>)(4\pi\tau)^{-\frac{n}{2}}e^{-f}dV \\
&\quad+2\tau\int_M[(\nabla_i\nabla_jf)R_{ij}
-\nabla_if\nabla_jfR_{ij}](4\pi\tau)^{-\frac{n}{2}}e^{-f}dV\\
&\quad+2\tau\int_M\(-\frac{1}{2\tau}g_{ij}\)(\nabla_i\nabla_jf)
(4\pi\tau)^{-\frac{n}{2}}e^{-f}dV\\
& =-2\tau\int_M[(\nabla_if\nabla_jf-\nabla_i\nabla_jf)
\nabla_i\nabla_jf-R_{ij}\nabla_if\nabla_jf\\
&\qquad-\nabla_i\nabla_j
f\nabla_if\nabla_jf](4\pi\tau)^{-\frac{n}{2}}e^{-f}dV\\
&\quad +2\tau\int_M[(\nabla_i\nabla_jf)R_{ij}
-\nabla_if\nabla_jfR_{ij}](4\pi\tau)^{-\frac{n}{2}}e^{-f}dV\\
&\quad +2\tau\int_M\(-\frac{1}{2\tau}g_{ij}\)(\nabla_i\nabla_jf)
(4\pi\tau)^{-\frac{n}{2}}e^{-f}dV\\
& =2\tau\int_M(\nabla_i\nabla_jf)\(\nabla_i\nabla_jf+R_{ij}
-\frac{1}{2\tau}g_{ij}\)(4\pi\tau)^{-\frac{n}{2}}e^{-f}dV.
\end{align*}
} Also the third term on the RHS of (1.5.11) is
\begin{align*}
& \int_M-\(R+|\nabla f|^2-\frac{n}{2\tau}\)
(4\pi\tau)^{-\frac{n}{2}}e^{-f}dV\\
&  = \int_M-\(R+\Delta f-\frac{n}{2\tau}\)
(4\pi\tau)^{-\frac{n}{2}}e^{-f}dV\\
& =2\tau\int_M\(\frac{-1}{2\tau}g_{ij}\)\(R_{ij}
+\nabla_i\nabla_jf-\frac{1}{2\tau}g_{ij}\)
(4\pi\tau)^{-\frac{n}{2}}e^{-f}dV.
\end{align*}
Therefore, by combining the above identities, we obtain
$$
\frac{d} {d t} \mathcal{W}(g_{ij}(t),f(t),\tau(t))
=2\tau\int_M\left|R_{ij}+\nabla_i\nabla_jf-\frac{1}{2\tau}g_{ij}\right|^2
(4\pi\tau)^{-\frac{n}{2}}e^{-f}dV.
$$

Finally, by using the computations in (1.5.5) and the evolution
equations of $f$ and $\tau$, we have
\begin{align*}
\frac{\partial}{\partial t}\((4\pi\tau)^{-\frac{n}{2}} e^{-f}dV\)
&= (4\pi\tau)^{-\frac{n}{2}}\left[\frac{\partial}{\partial t}(
e^{-f}dV)+\frac {n}{2\tau} e^{-f}dV\right]\\
&= -(4\pi\tau)^{-\frac{n}{2}}\Delta (e^{-f})dV.
\end{align*}
Hence
$$
\frac{d}{d t}\int_M(4\pi\tau)^{-\frac{n}{2}}e^{-f}dV
=-(4\pi\tau)^{-\frac{n}{2}}\int_M\Delta (e^{-f})dV=0.
$$
%%$$\eqno \#$$ \vskip 0.3cm
\end{proof}

Now we set \be
\mu(g_{ij},\tau)=\inf\left\{\mathcal{W}(g_{ij},f,\tau)\ |\  f\in
C^\infty(M),
\frac{1}{(4\pi\tau)^{n/2}}\int_M e^{-f}dV=1\right\}  %%\eqno(1.5.12)
\ee and
$$
\nu(g_{ij})=\inf\left\{\mathcal W(g,f,\tau)\ |\  f\in C^\infty(M),
\tau>0, \frac{1}{(4\pi\tau)^{n/2}}\int e^{-f}dV=1\right\}.
$$
Note that if we let $u=e^{-f/2}$, then the functional
$\mathcal{W}$ can be expressed as
$$
\mathcal{W}(g_{ij},f,\tau) =\int_M[\tau(Ru^2+4|\nabla
u|^2)-u^2\log u^2-nu^2] (4\pi\tau)^{-\frac{n}{2}}dV
$$
and the constraint $\int_M(4\pi\tau)^{-\frac{n}{2}}e^{-f}dV=1$
becomes $\int_Mu^2(4\pi\tau)^{-\frac{n}{2}}dV=1.$ Thus
$\mu(g_{ij},\tau)$ corresponds to the best constant of a
logarithmic Sobolev inequality. Since the nonquadratic term is
subcritical (in view of Sobolev exponent), it is rather
straightforward to show that
$$
\inf\!\bigg\{\!\int_M\![\tau(4|\nabla u|^2\!+Ru^2) -u^2\log
u^2\!-nu^2](4\pi\tau)^{-\frac{n}{2}}dV \Big|
\int_Mu^2(4\pi\tau)^{-\frac{n}{2}}dV\!=\!1\!\bigg\}
$$
is achieved by some nonnegative function $u\in H^1(M)$ which
satisfies the Euler-Lagrange equation
$$\tau (-4\Delta u+Ru)-2u\log u-nu=\mu(g_{ij},\tau)u
.$$ One can further show that $u$ is positive (see \cite{Ro}).
Then the standard regularity theory of elliptic PDEs shows that
$u$ is smooth. We refer the reader to Rothaus \cite{Ro} for more
details. It follows that $\mu(g_{ij},\tau)$ is achieved by a
minimizer $f$ satisfying the nonlinear equation \be
\tau (2\Delta f-|\nabla f|^2+R)+f-n=\mu(g_{ij},\tau).%\eqno(1.5.13)
\ee

%%\noindent{\bf Corollary 1.5.9}
% CZ-BASELINE source-line:2924 source-kind:corollary source-id:1.5.9
\begin{corollary}[mu monotonicity and shrinking breathers]
  \label{cor:cz-mu-monotonicity-and-shrinking-breathers}
  \dcref{Corollary 1.5.9}
  \group{Cao--Zhu Variational Structure and Dynamic Property}
  \source{cao-zhu}
\
\begin{itemize}
\item[(i)] $\mu(g_{ij}(t),\tau-t)$ is nondecreasing along the
Ricci flow; moveover, the monotonicity is strict unless we are on
a shrinking gradient soliton; \item[(ii)] A shrinking breather is
necessarily a shrinking gradient soliton.
\end{itemize}
\end{corollary}

%%\vskip 0.1cm \noindent{\bf Proof.}
\begin{proof}
  \uses{prop:cz-monotonicity-of-perelmans-w-functional}
Fix any time $t_0$, let $f_0$ be a minimizer of
$\mu(g_{ij}(t_0),\tau-t_0).$ Note that the backward heat equation
$$
\frac{\partial f}{\partial t}=-\Delta f+|\nabla
f|^2-R+\frac{n}{2\tau}
$$
is equivalent to the linear equation
$$
\frac{\partial }{\partial
t}((4\pi\tau)^{-\frac{n}{2}}e^{-f})=-\Delta
((4\pi\tau)^{-\frac{n}{2}}e^{-f})+R((4\pi\tau)^{-\frac{n}{2}}e^{-f}).
$$

Thus we can solve the backward heat equation of $f$ with
$f|_{t=t_0}=f_0$ to obtain $f(t)$, $t\leq t_0$, with
$\int_M(4\pi\tau)^{-\frac{n}{2}}e^{-f(t)}dV=1.$ It then follows
from Proposition 1.5.8 that
\begin{align*}
\mu(g_{ij}(t),\tau-t)
& \leq \mathcal{W}(g_{ij}(t),f(t),\tau-t)\\
& \leq \mathcal{W}(g_{ij}(t_0),f(t_0),\tau-t_0)\\
& = \mu(g_{ij}(t_0),\tau-t_0)
\end{align*}
for $t\leq t_0$ and the second inequality is strict unless we are
on a shrinking gradient soliton. This proves statement (i).

Consider a shrinking breather on $[t_1,t_2]$ with $\alpha
g_{ij}(t_1)$ and $g_{ij}(t_2)$ differ only by a diffeomorphism for
some $\alpha < 1.$ Recall that the functional $\mathcal{W}$ is
invariant under simultaneous scaling of $\tau$ and $g_{ij}$ and
invariant under diffeomorphism. Then for $\tau >0$ to be
determined,
$$
\mu(g_{ij}(t_1),\tau-t_1)=\mu(\alpha
g_{ij}(t_1),\alpha(\tau-t_1))=\mu(g_{ij}(t_2),\alpha(\tau-t_1))
$$
and by the monotonicity of $\mu(g_{ij}(t),\tau-t),$
$$
\mu(g_{ij}(t_1),\tau-t_1)\leq\mu(g_{ij}(t_2),\tau-t_2).
$$
Now take $\tau>0$ such that
$$
\alpha(\tau-t_1)=\tau-t_2,
$$
i.e.,
$$
\tau=\frac{t_2-\alpha t_1}{1-\alpha}.
$$

\noindent This shows the equality holds in the monotonicity of
$\mu(g_{ij}(t),\tau-t).$ So the shrinking breather must be a
shrinking gradient soliton.
%%$$\eqno \#$$ \vskip 0.3cm
\end{proof}

Finally, we remark that Hamilton, Ilmanen and the first author
\cite{Cao04} have obtained the second variation formulas for both
$\lambda$-energy and $\nu$-energy. We refer the reader to their
paper \cite{Cao04}  for more details and related stability
questions.

%%\bigskip
